% archon:covers Run202608192034.lean
% Historical Archon track only. The canonical Anderson/ proof is audited
% separately by StatementAudit.lean and StrictAxiomAudit.lean.
\chapter{The Anderson counterexample}

\section{The complete local node}

\begin{definition}[The node ring]
  \label{def:node_ring}
  \lean{Run202608192034.TODO.nodeRing}
  Let
  \[
    S=\mathbb C[[x,y,z]],\qquad
    T=S/(x^2-yz),
  \]
  and let \(\mathfrak M=(x,y,z)T\), where the letters also denote their
  residue classes.
\end{definition}

\begin{proof}
  This is a definition.
\end{proof}

\begin{definition}[The distinguished prime candidate]
  \label{def:node_prime}
  \lean{Run202608192034.TODO.nodePrime}
  \uses{def:node_ring}
  In \(T\), set \(Q=(x,y)T\).
\end{definition}

\begin{proof}
  This is a definition.
\end{proof}

\begin{lemma}[Normal form in the node ring]
  \label{lem:node_normal_form}
  \lean{Run202608192034.TODO.node_normal_form}
  \uses{def:node_ring}
  Every element of \(S\) has a unique expression
  \[
    A(y,z)+xB(y,z)+(x^2-yz)H(x,y,z)
  \]
  with \(A,B\in\mathbb C[[y,z]]\) and \(H\in S\).
\end{lemma}

\begin{proof}
  Regard \(x^2-yz\) as a monic polynomial in \(x\) over the complete
  coefficient ring \(\mathbb C[[y,z]]\).  Formal division by this monic
  polynomial gives a remainder of degree below two, hence the displayed
  expression.  If two such remainders differ by a multiple of the monic
  quadratic, comparing the \(x\)-degree below two forces both coefficients
  of the difference to vanish, proving uniqueness.
\end{proof}

\begin{lemma}[The node ring is a domain]
  \label{lem:node_domain}
  \lean{Run202608192034.TODO.nodeRing_isDomain}
  \uses{def:node_ring,lem:node_normal_form}
  The ring \(T\) is an integral domain.
\end{lemma}

\begin{proof}
  Map \(S\) continuously to \(\mathbb C[[u,v]]\) by
  \(x\mapsto uv\), \(y\mapsto u^2\), and \(z\mapsto v^2\).
  The relation \(x^2-yz\) maps to zero.  By
  \cref{lem:node_normal_form}, an element of the kernel has remainder
  \(A(y,z)+xB(y,z)\), whose image is
  \(A(u^2,v^2)+uvB(u^2,v^2)\).  The first summand contains only monomials
  with both exponents even and the second only monomials with both exponents
  odd.  They cannot cancel, so both vanish.  Thus the kernel is exactly
  \((x^2-yz)\), embedding \(T\) into the domain \(\mathbb C[[u,v]]\).
\end{proof}

\begin{lemma}[Completeness, dimension, and depth]
  \label{lem:node_complete_cm_dim}
  \lean{Run202608192034.TODO.node_complete_cm_dim}
  \uses{def:node_ring,lem:node_domain}
  The local ring \((T,\mathfrak M)\) is complete, Noetherian,
  two-dimensional, and Cohen--Macaulay of depth two.
\end{lemma}

\begin{proof}
  The power-series ring \(S\) is a complete regular local ring of dimension
  three.  By \cref{lem:node_domain}, the nonzero element \(x^2-yz\) is a
  nonzerodivisor.  A quotient of a complete Noetherian local ring is complete
  and Noetherian; quotienting a Cohen--Macaulay ring by a nonzerodivisor
  lowers both dimension and depth by one.  Hence \(T\) is a complete
  Cohen--Macaulay local ring of dimension and depth two.
\end{proof}

\begin{lemma}[Cardinality of the node ring]
  \label{lem:node_cardinality}
  \lean{Run202608192034.TODO.node_cardinality}
  \uses{def:node_ring}
  One has
  \[
    |T|=|T/\mathfrak M|=|\mathbb C|.
  \]
\end{lemma}

\begin{proof}
  There are countably many monomials in three variables, so
  \(|S|=|\mathbb C|^{\aleph_0}\).  Since
  \(|\mathbb C|=2^{\aleph_0}\), cardinal exponentiation gives
  \((2^{\aleph_0})^{\aleph_0}=2^{\aleph_0}\).  Constants inject into \(T\)
  and \(T\) is a quotient of \(S\), so Cantor--Schroeder--Bernstein gives
  \(|T|=|\mathbb C|\).  Finally \(T/\mathfrak M\cong\mathbb C\).
\end{proof}

\begin{lemma}[The distinguished ideal is prime of height one]
  \label{lem:node_prime_height}
  \lean{Run202608192034.TODO.nodePrime_prime_height}
  \uses{def:node_ring,def:node_prime,lem:node_domain,lem:node_complete_cm_dim}
  The ideal \(Q\) is a nonzero, nonmaximal prime ideal of \(T\) of height one.
\end{lemma}

\begin{proof}
  Setting \(x=y=0\) identifies \(T/Q\) with \(\mathbb C[[z]]\), so \(Q\)
  is prime.  It is nonzero because \(x\neq0\) in the domain \(T\), and it is
  nonmaximal because \(z\notin Q\).  Thus
  \((0)\subsetneq Q\subsetneq\mathfrak M\).  Since \(T\) has dimension two,
  the height of the intermediate prime is exactly one.
\end{proof}

\begin{lemma}[The distinguished prime is not principal]
  \label{lem:node_prime_nonprincipal}
  \lean{Run202608192034.TODO.nodePrime_not_principal}
  \uses{def:node_ring,def:node_prime,lem:node_prime_height}
  The ideal \(Q\) is not principal.
\end{lemma}

\begin{proof}
  The defining equation lies in the square of the maximal ideal, so the
  residue classes of \(x,y,z\) form a basis of
  \(\mathfrak M/\mathfrak M^2\) over \(\mathbb C\).  Since
  \(\mathfrak M Q\subseteq\mathfrak M^2\), the classes of \(x\) and \(y\)
  in \(Q/\mathfrak M Q\) are linearly independent.  Nakayama's lemma
  therefore shows that \(Q\) needs at least two generators, whereas a
  principal ideal needs at most one.
\end{proof}

\begin{lemma}[Noetherianity transports across a ring equivalence]
  \label{lem:is_noetherian_ring_of_ring_equiv_source}
  \lean{Run202608192034.TODO.isNoetherianRing_of_ringEquiv_source}
  If \(R\cong S\) as rings and \(R\) is Noetherian, then \(S\) is
  Noetherian.
\end{lemma}

\begin{proof}
  Rewrite Noetherianity as the module condition `IsNoetherian R R`.
  The ring equivalence gives a semilinear equivalence from \(R\) to \(S\),
  and Mathlib's transfer theorem for Noetherian modules along bijective
  semilinear maps transports the condition.
\end{proof}

\begin{lemma}[Complete domain package]
  \label{lem:complete_domain_choice}
  \lean{Run202608192034.TODO.completeDomainChoice}
  \uses{lem:node_domain,lem:node_complete_cm_dim,lem:node_cardinality,lem:node_prime_height,lem:node_prime_nonprincipal}
  The pair \((T,\mathfrak M)\) is a complete two-dimensional
  Cohen--Macaulay local domain with \(|T|=|T/\mathfrak M|\), and \(Q\) is a
  nonprincipal height-one prime of \(T\).
\end{lemma}

\begin{proof}
  The preceding paragraphs give the source-level decomposition.  In Lean v4.29
  the finite-variable formal-power-series division, Noetherianity, and dimension
  APIs needed to internalize that decomposition are unavailable, so the theorem
  invokes the visible standard-facts boundary
  `nodeRing_standard_facts_external`.  All downstream projections are checked;
  the exact trust boundary and internalization plan are recorded in
  `TRUST_BOUNDARY.md`.
\end{proof}

\section{The Jensen construction}

\begin{definition}[N-subring]
  \label{def:n_subring}
  \lean{Run202608192034.TODO.NSubring}
  % SOURCE: references/jensen_2006_fulltext.txt, lines 188--194.
  \textit{Source: D.~Jensen, Definition~1,
  references/jensen\_2006\_fulltext.txt, lines 188--194,
  DOI 10.1080/00927870500346321.}
  Let \((T,\mathfrak M)\) be complete local.  An N-subring is a quasi-local
  UFD \(R\subseteq T\) such that its cardinality is at most
  \(\max(\aleph_0,|T/\mathfrak M|)\), equality being allowed only when the
  residue field is countable; every associated prime of \(T\) contracts to
  zero in \(R\); and, for every regular \(t\in T\) and every
  \(P\in\operatorname{Ass}(T/tT)\), the prime \(P\cap R\) has height at most
  one.
\end{definition}

\begin{proof}
  This is the auxiliary definition used to preserve unique factorization
  during the transfinite construction.
\end{proof}

\begin{lemma}[Cardinal prime avoidance]
  \label{lem:cardinal_prime_avoidance}
  \lean{Run202608192034.TODO.cardinal_prime_avoidance}
  % SOURCE: references/loepp_1997_fulltext.txt, lines 153--167; references/heitmann_1993_fulltext.txt, lines 98--138.
  \textit{Source: S.~Loepp, Lemmas~2--4,
  references/loepp\_1997\_fulltext.txt, lines 153--167,
  DOI 10.1006/jabr.1997.6768; R.~C. Heitmann, Lemmas~2--3,
  references/heitmann\_1993\_fulltext.txt, lines 98--138,
  DOI 10.1090/S0002-9947-1993-1102888-9.}
  Let \((T,\mathfrak M)\) be a complete local ring, let \(C\) be a family of
  nonmaximal prime ideals, and for each \(P\in C\) let \(D_P\) be a set of
  forbidden coset representatives in \(T/P\).  If an ideal \(J\) is
  contained in no member of \(C\) and the set of pairs
  \(\{(P,d):P\in C,\ d\in D_P\}\) has cardinality strictly below
  \(|T/\mathfrak M|\), then \(J\) contains an element outside every coset
  \(P+d\).
\end{lemma}

\begin{proof}
  As \(T\) is Noetherian, write \(J=(x_1,\ldots,x_n)\).  For every controlled
  prime \(P\), let \(i(P)\) be the least index with \(x_{i(P)}\notin P\).
  Choose residue classes \(a_n,\ldots,a_1\in T/\mathfrak M\) in descending
  order.  When \(a_i\) is chosen, consider the forbidden pairs with
  \(i(P)=i\).  The already fixed terms have larger index, while every term
  of smaller index lies in \(P\).  A given pair \((P,d)\) rules out at most
  one residue class for \(a_i\): two distinct choices would differ by a unit
  times \(x_i\), forcing \(x_i\in P\).  Fewer than
  \(|T/\mathfrak M|\) classes are forbidden, so a permissible class exists.
  After lifting the \(a_i\) to \(T\), the element
  \(\sum_i a_i x_i\in J\) avoids every prescribed coset.
\end{proof}

\begin{lemma}[The Jensen residue field is uncountable]
  \label{lem:jensen_residue_uncountable}
  \lean{Run202608192034.TODO.jensen_residueField_uncountable}
  Let \((T,\mathfrak M)\) be a complete local ring of positive depth.  If
  \(|T|=|T/\mathfrak M|\), then the residue field \(T/\mathfrak M\) is
  uncountable.
\end{lemma}

\begin{proof}
  Choose a regular element \(x\in\mathfrak M\).  After choosing one lift in
  \(T\) of every residue class, completeness assigns to each sequence
  \((a_n)_{n\geq0}\) in the residue field the convergent series
  \(\sum_{n\geq0}\widetilde a_nx^n\).  This assignment is injective.  Indeed,
  at the first index where two sequences differ, their difference is
  \(x^n(u+xv)\), where \(u\) is a unit; regularity of \(x\) makes this
  nonzero.  Thus \(|T|\geq|T/\mathfrak M|^{\aleph_0}\).  A finite or countably
  infinite field has uncountably many sequences, contradicting the assumed
  equality.  Hence the residue field is uncountable.
\end{proof}

\begin{lemma}[Initial N-subring]
  \label{lem:initial_n_subring}
  \lean{Run202608192034.TODO.initialNSubring}
  \uses{def:n_subring}
  % SOURCE: references/jensen_2006_fulltext.txt, lines 245--293; references/loepp_1997_fulltext.txt, lines 997--1003.
  \textit{Source: D.~Jensen, proof of Theorem~2.2,
  references/jensen\_2006\_fulltext.txt, lines 245--293,
  DOI 10.1080/00927870500346321; S.~Loepp, proof of Theorem~16,
  references/loepp\_1997\_fulltext.txt, lines 997--1003,
  DOI 10.1006/jabr.1997.6768.}
  Let \((T,\mathfrak M)\) and \(G\) satisfy the hypotheses of the Jensen
  semilocal generic-fiber construction: in particular, every member of
  \(G\) meets the prime subring in zero, \(G\) contains
  \(\operatorname{Ass}T\), and \(G\) contains every prime \(P\) for which
  \(\operatorname{ht}P>\operatorname{depth}T_P=1\).  Then the prime subring
  of \(T\), localized at its intersection with \(\mathfrak M\), is an
  N-subring and has zero intersection with every prime in \(G\).
\end{lemma}

\begin{proof}
  The localized prime subring is a countable quasi-local UFD and avoids
  \(G\) by the prime-subring hypothesis.  Since \(G\) contains
  \(\operatorname{Ass}T\), associated primes contract to zero.  For regular
  \(t\) and \(P\in\operatorname{Ass}(T/tT)\), localization gives
  \(\operatorname{depth}T_P=1\).  If \(\operatorname{ht}P>1\), the hypotheses
  place \(P\) in \(G\), so its contraction is zero; otherwise its contraction
  has height at most one.  Thus the localized prime subring itself satisfies
  every N-subring condition and has the asserted contractions.
\end{proof}

\begin{lemma}[N-subring prime-hitting extension]
  \label{lem:n_subring_prime_extension}
  \lean{Run202608192034.TODO.nSubring_prime_extension}
  \uses{def:n_subring,lem:cardinal_prime_avoidance}
  % SOURCE: references/jensen_2006_fulltext.txt, lines 201--233.
  \textit{Source: D.~Jensen, Lemma~2.1,
  references/jensen\_2006\_fulltext.txt, lines 201--233,
  DOI 10.1080/00927870500346321.}
  Let \((T,\mathfrak M)\) be complete local of dimension at least two, and let
  \(G\subseteq\operatorname{Spec}T\) be downward closed, with finitely
  many nonmaximal maximal elements, containing \(\operatorname{Ass}T\) and
  every prime \(J\) with
  \(\operatorname{ht}J>\operatorname{depth}T_J=1\).  If an N-subring \(R\)
  contracts every member of \(G\) to zero and \(Q\) is a prime outside
  \(G\), then there is
  an N-subring \(S\) containing \(R\) such that \(Q\cap S\neq0\), every
  member of \(G\) still contracts to zero, the required cardinal bound is
  preserved, and prime elements of \(R\) remain prime in \(S\).
\end{lemma}

\begin{proof}
  Form the cardinality-controlled family \(C\) consisting of the maximal
  members of \(G\) and the associated primes of \(T/rT\) for nonzero
  \(r\in R\).  If \(Q\in C\), then \(Q\notin G\) forces
  \(Q\in\operatorname{Ass}(T/rT)\) for some nonzero \(r\in R\), and already
  \(0\neq r\in Q\cap R\), so take \(S=R\).  Otherwise \(Q\) is contained in
  no member of \(C\): downward closure handles maximal members of \(G\); and
  an associated prime of \(T/rT\) is either in \(G\), or has height one by
  the N-subring and depth hypotheses, in which case a strictly smaller prime
  is minimal, hence associated to \(T\) and therefore lies in \(G\).

  For each member of \(C\), forbid the cosets algebraic over the
  corresponding image of \(R\).  Their total cardinality obeys the N-subring
  bound, so \cref{lem:cardinal_prime_avoidance} supplies
  \(t\in Q\) transcendental modulo every forbidden prime.  Localize
  \(R[t]\) at its intersection with \(\mathfrak M\).  Transcendence gives
  zero contraction for every maximal member of \(G\), hence for all of
  \(G\); it also preserves the N-subring conditions and old prime elements,
  while \(t\) makes the new ring meet \(Q\) nontrivially.
\end{proof}

\begin{lemma}[N-subring ideal-solving extension]
  \label{lem:n_subring_ideal_extension}
  \lean{Run202608192034.TODO.nSubring_ideal_extension}
  \uses{def:n_subring,lem:cardinal_prime_avoidance}
  % SOURCE: references/heitmann_1993_fulltext.txt, lines 145--148 and 304--331; references/loepp_1997_fulltext.txt, lines 907--977.
  \textit{Source: R.~C. Heitmann, Lemma~4 and Lemma~7,
  references/heitmann\_1993\_fulltext.txt, lines 145--148 and 304--331,
  DOI 10.1090/S0002-9947-1993-1102888-9; S.~Loepp, Lemma~15,
  references/loepp\_1997\_fulltext.txt, lines 907--977,
  DOI 10.1006/jabr.1997.6768.}
  Under the same avoidance hypotheses, let \(I\) be a finitely generated
  ideal of an N-subring \(R\), and let \(c\in IT\cap R\).  There is an
  N-subring extension \(S\) which preserves avoidance and old prime elements
  and in which \(c\in IS\).  Likewise, any prescribed coset of
  \(T/\mathfrak M^2\) can be represented by an element of such an extension.
\end{lemma}

\begin{proof}
  Write \(c\) as a \(T\)-linear combination of fixed generators of \(I\).
  Adjoin the finitely many required coefficients one at a time, perturbing
  each by an element of a high power of \(\mathfrak M\).  At every step,
  \cref{lem:cardinal_prime_avoidance} chooses the perturbation so that its
  residue remains transcendental modulo all controlled primes.  The resulting
  localized polynomial extension remains an N-subring, realizes
  \(c\in IS\), and preserves all contractions and prime elements.  The same
  argument with a chosen representative solves a prescribed coset modulo
  \(\mathfrak M^2\).
\end{proof}

\begin{definition}[Jensen saturation chain]
  \label{def:jensen_saturation_chain}
  \lean{Run202608192034.TODO.jensenSaturationChain}
  \uses{lem:n_subring_prime_extension,lem:n_subring_ideal_extension}
  % SOURCE: references/loepp_1997_fulltext.txt, lines 802--1033; references/heitmann_1993_fulltext.txt, lines 264--331.
  \textit{Source: S.~Loepp, Lemmas~13--15 and Theorem~16,
  references/loepp\_1997\_fulltext.txt, lines 802--1033,
  DOI 10.1006/jabr.1997.6768; R.~C. Heitmann, Lemmas~6--7,
  references/heitmann\_1993\_fulltext.txt, lines 264--331,
  DOI 10.1090/S0002-9947-1993-1102888-9.}
  Fix well-orderings of the cosets of \(T/\mathfrak M^2\), the primes outside
  \(G\), and all pairs \((I,c)\) which arise with \(I\) finitely generated in
  a stage ring and \(c\in IT\) belonging to that stage.  A Jensen saturation
  chain is a continuous transfinite chain of N-subrings which applies the
  appropriate extension lemma cofinally often to every item in these
  well-orderings.
\end{definition}

\begin{proof}
  Transfinite recursion constructs the chain.  At successor stages use
  \cref{lem:n_subring_prime_extension} for the next prime and
  \cref{lem:n_subring_ideal_extension} for the next coset or ideal equation.
  At a limit stage take the directed union and localize at its intersection
  with \(\mathfrak M\).  The cardinal bounds and zero-contraction conditions
  are preserved under such unions, and the repeated bookkeeping ensures that
  every task appearing at any stage is handled later.
\end{proof}

\begin{lemma}[The saturated union is a UFD]
  \label{lem:jensen_union_ufd}
  \lean{Run202608192034.TODO.jensenUnion_isUFD}
  \uses{def:jensen_saturation_chain,def:n_subring}
  The union \(A\) of a Jensen saturation chain is a quasi-local UFD, and every
  prime in \(G\) has zero contraction to \(A\).
\end{lemma}

\begin{proof}
  Every nonzero nonunit of \(A\) lies in some stage, where it factors into
  prime elements because that stage is a UFD.  Successor extensions preserve
  those prime elements, so the same factorization is prime in every later
  stage and in the union.  Thus \(A\) is a UFD.  Quasi-locality follows
  because all stages use the maximal ideal induced by \(\mathfrak M\).
  Each stage avoids every prime in \(G\), and hence their union does too.
\end{proof}

\begin{lemma}[Completion criterion for the saturated union]
  \label{lem:jensen_completion_criterion}
  \lean{Run202608192034.TODO.jensen_completion_criterion}
  \uses{def:jensen_saturation_chain,lem:jensen_union_ufd}
  % SOURCE: references/loepp_1997_fulltext.txt, lines 141--144 and 1024--1033; references/heitmann_1993_fulltext.txt, lines 332--353.
  \textit{Source: S.~Loepp, Proposition~1 and proof of Theorem~16,
  references/loepp\_1997\_fulltext.txt, lines 141--144 and 1024--1033,
  DOI 10.1006/jabr.1997.6768; R.~C. Heitmann, Theorem~8,
  references/heitmann\_1993\_fulltext.txt, lines 332--353,
  DOI 10.1090/S0002-9947-1993-1102888-9.}
  Let \(A\subseteq T\) be the saturated union.  If
  \(A\to T/\mathfrak M^2\) is onto and \(IT\cap A=I\) for every finitely
  generated ideal \(I\) of \(A\), then \(A\) is Noetherian and local, and the
  natural map from its completion to \(T\) is an isomorphism.
\end{lemma}

\begin{proof}
  Surjectivity modulo \(\mathfrak M^2\) first identifies the maximal ideal of
  \(A\) and its cotangent space with those induced from \(T\).  The contraction
  identities imply exactness after extending every finitely generated ideal
  to \(T\); in particular, finitely generated subideals which generate the
  same extended ideal already generate the original ideal.  For an arbitrary
  ideal \(J\subseteq A\), Noetherianity of \(T\) supplies
  \(j_1,\ldots,j_r\in J\) whose extensions generate \(JT\).  Put
  \(I=(j_1,\ldots,j_r)A\).  The contraction hypothesis gives
  \(J\subseteq JT\cap A=IT\cap A=I\subseteq J\), hence \(J=I\).  Thus every
  ideal of \(A\) is finitely generated and \(A\) is Noetherian.
  Inductively, the same identities identify
  \(A/\mathfrak m_A^n\) with \(T/\mathfrak M^n\) for all \(n\).  Taking inverse
  limits gives \(\widehat A\cong T\).
\end{proof}

\begin{theorem}[Jensen semilocal generic-fiber construction]
  \label{thm:jensen_semilocal_generic_fiber}
  \lean{Run202608192034.TODO.jensen_semilocal_genericFiber}
  \uses{def:n_subring,lem:initial_n_subring,lem:n_subring_prime_extension,lem:n_subring_ideal_extension,def:jensen_saturation_chain,lem:jensen_union_ufd,lem:jensen_completion_criterion,def:generic_formal_fiber}
  % SOURCE: references/jensen_2006_fulltext.txt, lines 245--293.
  \textit{Source: D.~Jensen, Theorem~2.2,
  references/jensen\_2006\_fulltext.txt, lines 245--293,
  DOI 10.1080/00927870500346321.}
  Let \((T,\mathfrak M)\) be complete local with
  \(|T/\mathfrak M|=|T|\).  Let \(G\) be a nonempty downward-closed subset of
  \(\operatorname{Spec}T\) with finitely many maximal elements.  Assume that
  \(\mathfrak M\notin G\), that \(G\) contains \(\operatorname{Ass}T\), that
  every member of \(G\) meets the prime subring only in zero, and that every
  prime \(J\) with
  \(\operatorname{ht}J>\operatorname{depth}T_J=1\) belongs to \(G\).
  Then there is a Noetherian local UFD \(A\) whose completion is \(T\) and
  whose generic formal fiber consists exactly of \(G\).
\end{theorem}

\begin{proof}
  Start with the localized prime subring supplied by
  \cref{lem:initial_n_subring}.  Build the saturation chain of
  \cref{def:jensen_saturation_chain}.  Its union is a UFD avoiding every
  member of \(G\) by \cref{lem:jensen_union_ufd}.  The cofinal coset and
  ideal-solving tasks give the hypotheses of
  \cref{lem:jensen_completion_criterion}, so the union is Noetherian local
  and completes to \(T\).  Each prime outside \(G\) is assigned a
  prime-hitting stage and therefore has nonzero contraction, while every
  prime in \(G\) has zero contraction.  By the prime description of
  \cref{def:generic_formal_fiber}, the generic formal fiber is exactly \(G\).
\end{proof}

\begin{corollary}[Jensen local generic-fiber criterion]
  \label{cor:jensen_local_generic_fiber}
  \lean{Run202608192034.TODO.jensen_local_genericFiber}
  \uses{thm:jensen_semilocal_generic_fiber,lem:jensen_residue_uncountable}
  % SOURCE: references/jensen_2006_fulltext.txt, lines 336--342.
  \textit{Source: D.~Jensen, Corollary~2.4,
  references/jensen\_2006\_fulltext.txt, lines 336--342,
  DOI 10.1080/00927870500346321.}
  Let \((T,\mathfrak M)\) be complete local with
  \(|T/\mathfrak M|=|T|\), and let \(P\) be prime.  If \(T\) has depth at
  least two, \(P\) is nonmaximal and contains every associated prime of
  \(T\), \(P\) meets the prime subring only in zero, and every prime \(J\)
  satisfying \(\operatorname{ht}J>\operatorname{depth}T_J=1\) is contained
  in \(P\), then there is a local UFD \(A\) completing to \(T\) whose generic
  formal fiber is local with maximal prime \(P\).
\end{corollary}

\begin{proof}
  Apply \cref{thm:jensen_semilocal_generic_fiber} to
  \(G=\{Q\in\operatorname{Spec}T:Q\subseteq P\}\).  This set is nonempty,
  downward closed, and has the unique maximal element \(P\).  The stated
  hypotheses are precisely the remaining hypotheses of that theorem.
\end{proof}

\begin{lemma}[The node satisfies Jensen's hypotheses]
  \label{lem:node_jensen_hypotheses}
  \lean{Run202608192034.TODO.node_jensen_hypotheses}
  \uses{lem:complete_domain_choice,cor:jensen_local_generic_fiber}
  The ring \(T\) and the prime \(P=(0)\) satisfy all hypotheses of
  \cref{cor:jensen_local_generic_fiber}.
\end{lemma}

\begin{proof}
  Completeness, the cardinal equality, and depth two are part of
  \cref{lem:complete_domain_choice}.  Since \(T\) is a domain,
  \(\operatorname{Ass}T=\{(0)\}\); the zero prime is nonmaximal and has zero
  intersection with the prime subring.  Every localization of the
  Cohen--Macaulay ring \(T\) is Cohen--Macaulay.  Thus for every prime \(J\),
  \(\operatorname{depth}T_J=\dim T_J=\operatorname{ht}J\), so no prime
  satisfies \(\operatorname{ht}J>\operatorname{depth}T_J=1\).
\end{proof}

\begin{definition}[Jensen initial N-subring seed data]
  \label{def:jensen_n_subring_seed_data}
  \lean{Run202608192034.TODO.JensenNSubringSeedData}
  \uses{def:n_subring}
  Record the initial data for Jensen's construction: the chosen generic-fiber
  set \(G=\{(0)\}\), an initial \(N\)-subring \(R_0\subseteq T\), and the
  initial zero-contraction condition for primes in \(G\).
\end{definition}

\begin{proof}
  This is a data structure matching the first setup step of Jensen's proof.
\end{proof}

\begin{definition}[Jensen saturation data]
  \label{def:jensen_saturation_data}
  \lean{Run202608192034.TODO.JensenSaturationData}
  \uses{def:jensen_n_subring_seed_data,def:jensen_saturation_chain}
  Record the saturation stage: a chain of \(N\)-subrings, its union \(A\), the
  preservation of zero contraction on \(G\), and the prime-hitting condition for
  primes outside \(G\).
\end{definition}

\begin{proof}
  This packages the bookkeeping performed by Jensen's recursive construction.
\end{proof}

\begin{definition}[Jensen UFD output data]
  \label{def:jensen_ufd_output_data}
  \lean{Run202608192034.TODO.JensenUFDOutputData}
  \uses{def:jensen_saturation_data}
  Record the conclusion that the saturated union \(A\) is a UFD.
\end{definition}

\begin{proof}
  This is the UFD output of the saturation argument.
\end{proof}

\begin{definition}[Jensen completion criterion output data]
  \label{def:jensen_completion_criterion_output_data}
  \lean{Run202608192034.TODO.JensenCompletionCriterionOutputData}
  \uses{def:jensen_saturation_data,lem:jensen_completion_criterion}
  Record the completion-criterion output: \(A\) is Noetherian and local, the
  selected ideal is the maximal ideal, \(\widehat A\cong T\), the map
  \(A\to T\) is compatible with this completion equivalence, and
  \(\dim A=2\).
\end{definition}

\begin{proof}
  This packages the conclusion obtained after the saturation chain solves the
  required coset and finitely generated ideal equations.
\end{proof}

\begin{definition}[Jensen generic formal fiber output data]
  \label{def:jensen_generic_fiber_output_data}
  \lean{Run202608192034.TODO.JensenGenericFiberOutputData}
  \uses{def:jensen_saturation_data,def:jensen_completion_criterion_output_data,def:generic_formal_fiber}
  Record the final generic-formal-fiber output for the selected \(A\): the zero
  prime contracts to zero and every nonzero prime of \(T\) contracts
  nontrivially.
\end{definition}

\begin{proof}
  This packages only the prime description supplied by Jensen.  The weak
  quasi-completeness criterion is added later from Anderson/Farley, rather than
  being attributed to Jensen's construction.
\end{proof}

\begin{lemma}[Initial N-subring seed source]
  \label{lem:jensen_n_subring_seed_source}
  \lean{Run202608192034.TODO.jensen_nSubring_seed_source}
  \uses{def:jensen_published_construction_data}
  A certificate for Jensen's published construction contains its initial
  \(N\)-subring seed for \(G=\{(0)\}\).
\end{lemma}

\begin{proof}
  This is a checked field projection.  The construction hidden by the
  published-source boundary starts from Jensen's localized prime subring and
  applies the initial \(N\)-subring lemma.
\end{proof}

\begin{definition}[Jensen saturation task set]
  \label{def:jensen_saturation_task_enumeration}
  \lean{Run202608192034.TODO.JensenSaturationTaskEnumeration}
  \uses{def:jensen_n_subring_seed_data}
  Record the bookkeeping device that collects the prime-hitting tasks and
  ideal-solving tasks used by Jensen's recursive saturation construction.
\end{definition}

\begin{proof}
  This is data: the prime task set covers the primes outside \(G\), while the
  ideal tasks record the algebraic equations that must be solved along the
  chain.
\end{proof}

\begin{definition}[Jensen recursive saturation chain data]
  \label{def:jensen_saturation_recursive_chain_data}
  \lean{Run202608192034.TODO.JensenSaturationRecursiveChainData}
  \uses{def:jensen_saturation_task_enumeration,def:jensen_saturation_chain,def:jensen_n_subring_seed_data}
  Record the recursive chain produced from the task enumeration: it starts
  above the initial \(N\)-subring, preserves avoidance at each finite stage,
  and eventually hits every prime outside \(G\).
\end{definition}

\begin{proof}
  This packages the output expected from the successor-stage construction using
  prime extension and ideal extension.
\end{proof}

\begin{definition}[Jensen saturation union data]
  \label{def:jensen_saturation_union_data}
  \lean{Run202608192034.TODO.JensenSaturationUnionData}
  \uses{def:jensen_saturation_recursive_chain_data}
  Record the union ring \(A\) of the saturation chain, together with the facts
  that every finite stage maps into \(A\), \(G\)-avoidance survives the union,
  and primes outside \(G\) still have nonzero contraction.
\end{definition}

\begin{proof}
  This is the limit-stage output of Jensen's construction.
\end{proof}

\begin{lemma}[Saturation task set source]
  \label{lem:jensen_saturation_task_enumeration_source}
  \lean{Run202608192034.TODO.jensen_saturation_task_enumeration_source}
  \uses{def:jensen_saturation_task_enumeration,def:jensen_n_subring_seed_data}
  For a fixed seed, collect every prime outside \(G\) as a prime-hitting task
  and every seed ideal as an ideal-solving task.
\end{lemma}

\begin{proof}
  Lean constructs these two sets directly and checks the prime-coverage field.
\end{proof}

\begin{lemma}[Saturation recursive chain source]
  \label{lem:jensen_saturation_recursive_chain_source}
  \lean{Run202608192034.TODO.jensen_saturation_recursive_chain_source}
  \uses{def:jensen_published_construction_data,def:jensen_saturation_recursive_chain_data}
  A certificate for Jensen's published construction contains the recursive
  saturation chain, with stagewise avoidance and eventual prime hitting.
\end{lemma}

\begin{proof}
  This is a checked field projection from the published construction
  certificate.  Internally, Jensen obtains the field by iterating the prime-
  and ideal-extension lemmas.
\end{proof}

\begin{lemma}[Saturation union source]
  \label{lem:jensen_saturation_union_source}
  \lean{Run202608192034.TODO.jensen_saturation_union_source}
  \uses{def:jensen_saturation_recursive_chain_data,def:jensen_saturation_union_data}
  The union of the recursive saturation chain is a subring \(A\) preserving
  \(G\)-avoidance and the eventual prime-hitting property.
\end{lemma}

\begin{proof}
  Lean sets \(A=\bigvee_n R_n\).  Monotonicity of the chain makes the family
  directed, so membership in the supremum comes from a finite stage.  The
  stagewise zero-contraction property then proves avoidance, while an eventual
  hitting witness embeds into the supremum and proves nonzero contraction.
\end{proof}

\begin{lemma}[Recursive chain starts]
  \label{lem:jensen_saturation_recursive_chain_starts}
  \lean{Run202608192034.TODO.JensenSaturationRecursiveChainData.starts}
  \uses{def:jensen_saturation_recursive_chain_data}
  The recursive saturation chain starts above the initial \(N\)-subring.
\end{lemma}

\begin{proof}
  This is the corresponding field projection from the recursive-chain data.
\end{proof}

\begin{lemma}[Union contains chain]
  \label{lem:jensen_saturation_union_chain_le}
  \lean{Run202608192034.TODO.JensenSaturationUnionData.chain_le}
  \uses{def:jensen_saturation_union_data}
  Every finite stage of the saturation chain is contained in the union ring.
\end{lemma}

\begin{proof}
  This is the corresponding field projection from the union data.
\end{proof}

\begin{lemma}[Union preserves avoidance]
  \label{lem:jensen_saturation_union_avoid}
  \lean{Run202608192034.TODO.JensenSaturationUnionData.avoid}
  \uses{def:jensen_saturation_union_data}
  The union ring still has zero contraction for all primes in \(G\).
\end{lemma}

\begin{proof}
  This is the corresponding field projection from the union data.
\end{proof}

\begin{lemma}[Union hits primes outside \(G\)]
  \label{lem:jensen_saturation_union_hits_outside_g}
  \lean{Run202608192034.TODO.JensenSaturationUnionData.hitsOutsideG}
  \uses{def:jensen_saturation_union_data}
  Every prime outside \(G\) has nonzero contraction to the union ring.
\end{lemma}

\begin{proof}
  This is the corresponding field projection from the union data.
\end{proof}

\begin{definition}[Union data as saturation data]
  \label{def:jensen_saturation_union_to_saturation_data}
  \lean{Run202608192034.TODO.JensenSaturationUnionData.toSaturationData}
  \uses{def:jensen_saturation_union_data,def:jensen_saturation_data,lem:jensen_saturation_recursive_chain_starts,lem:jensen_saturation_union_chain_le,lem:jensen_saturation_union_avoid,lem:jensen_saturation_union_hits_outside_g}
  Convert the recursive-chain and union bookkeeping package into the
  saturation data used by the rest of the proof.
\end{definition}

\begin{proof}
  Assemble the fields of `JensenSaturationData` from the checked projections.
\end{proof}

\begin{definition}[Saturation data selected from a recursive chain]
  \label{def:jensen_saturation_data_of_recursive_chain}
  \lean{Run202608192034.TODO.jensenSaturationDataOfRecursiveChain}
  \uses{lem:jensen_saturation_union_source,def:jensen_saturation_union_to_saturation_data}
  Select the checked directed union of a recursive chain and expose it as the
  saturation-data package used by the later dependent outputs.
\end{definition}

\begin{proof}
  Apply classical choice to \cref{lem:jensen_saturation_union_source}, then
  apply the checked conversion to saturation data.
\end{proof}

\begin{definition}[Published Jensen construction certificate]
  \label{def:jensen_published_construction_data}
  \lean{Run202608192034.TODO.JensenPublishedConstructionData}
  \uses{def:jensen_n_subring_seed_data,def:jensen_saturation_task_enumeration,def:jensen_saturation_recursive_chain_data,def:jensen_saturation_data_of_recursive_chain,def:jensen_ufd_output_data,def:jensen_completion_criterion_output_data,def:jensen_generic_fiber_output_data}
  Package one coherent dependent trace of Jensen's construction: its seed,
  task enumeration, recursive chain, saturated union, UFD output, completion
  output, and generic-formal-fiber output.
\end{definition}

\begin{proof}
  This is a dependent record.  Lean assumes its nonemptiness through the
  explicitly named published-source boundary
  \texttt{jensen\_corollary\_2\_4\_construction\_external}, which is Jensen,
  Corollary~2.4 specialized to the node ring and \(P=(0)\).  See
  \texttt{TRUST\_BOUNDARY.md} for the exact audit.
\end{proof}

\begin{lemma}[Saturation source]
  \label{lem:jensen_saturation_source}
  \lean{Run202608192034.TODO.jensen_saturation_source}
  \uses{def:jensen_published_construction_data,def:jensen_saturation_data_of_recursive_chain}
  A published construction certificate determines the saturated union \(A\),
  preserving \(G\)-avoidance and hitting every prime outside \(G\).
\end{lemma}

\begin{proof}
  Apply the checked saturation-data constructor to the certificate's recursive
  chain.
\end{proof}

\begin{lemma}[UFD output source]
  \label{lem:jensen_ufd_output_source}
  \lean{Run202608192034.TODO.jensen_ufd_output_source}
  \uses{def:jensen_published_construction_data,def:jensen_ufd_output_data}
  The saturated union produced by Jensen's chain is a UFD.
\end{lemma}

\begin{proof}
  This is a checked field projection from the coherent Jensen certificate.
\end{proof}

\begin{lemma}[Completion output source]
  \label{lem:jensen_completion_output_source}
  \lean{Run202608192034.TODO.jensen_completion_output_source}
  \uses{def:jensen_published_construction_data,def:jensen_completion_criterion_output_data}
  The saturated union satisfies the completion criterion, so its completion is
  \(T\), it is Noetherian local, and it has dimension two.
\end{lemma}

\begin{proof}
  This is a checked field projection from the coherent Jensen certificate.
  Jensen obtains it from the coset-surjectivity and ideal-contraction
  equations built into the saturation process.
\end{proof}

\begin{lemma}[Generic formal fiber output source]
  \label{lem:jensen_generic_formal_fiber_output_source}
  \lean{Run202608192034.TODO.jensen_generic_formal_fiber_output_source}
  \uses{def:jensen_published_construction_data,def:jensen_generic_fiber_output_data}
  The selected source has trivial generic formal fiber: zero contracts to zero
  and every nonzero prime of \(T\) contracts nontrivially.
\end{lemma}

\begin{proof}
  This is a checked field projection from the coherent Jensen certificate.
  In the source proof, \(G=\{(0)\}\), so avoidance and hitting identify the
  generic formal fiber exactly.
\end{proof}

\begin{definition}[Jensen selected source]
  \label{def:jensen_selected_source}
  \lean{Run202608192034.TODO.JensenSelectedSource}
  \uses{def:jensen_completion_witness,def:jensen_saturation_data,def:jensen_completion_criterion_output_data,def:jensen_generic_fiber_output_data}
  Package the actual source ring selected by Jensen's construction: a
  Noetherian local domain \(A\), its ideal \(\mathfrak m\), the map
  \(A\to T\), the completion witness \(\widehat A\cong T\), and the
  generic-formal-fiber contraction condition.
\end{definition}

\begin{proof}
  This is a data package.  It separates the construction of Jensen's selected
  source ring from the later projections of its UFD, completion, dimension, and
  contraction properties.
\end{proof}

\begin{lemma}[Jensen selected source exists]
  \label{lem:jensen_selected_source_exists_source}
  \lean{Run202608192034.TODO.jensenSelectedSource_exists_source}
  \uses{def:jensen_published_construction_data,def:jensen_selected_source,def:jensen_saturation_data_of_recursive_chain,lem:farley_completion_criterion,lem:adic_completion_prime_contraction_condition_iff_target}
  Jensen's published construction certificate assembles into the selected
  source package used downstream.
\end{lemma}

\begin{proof}
  Choose the certificate supplied by the explicit Corollary~2.4 boundary.
  Anderson's Corollary~2(1) gives the weak-completeness criterion on the actual
  adic completion, and the checked prime-contraction transport lemma carries it
  across \(\widehat A\cong T\).  Assemble this with the certificate's dependent
  UFD, completion, and generic-formal-fiber fields into
  `JensenCompletionWitness` and `JensenSelectedSource`.
\end{proof}

\begin{lemma}[Jensen selected ring exists]
  \label{lem:jensen_selected_ring_exists_source}
  \lean{Run202608192034.TODO.jensen_selected_ring_exists_source}
  \uses{lem:jensen_selected_source_exists_source}
  The selected source package contains a Noetherian local domain \(A\) and a map
  \(A\to T\).
\end{lemma}

\begin{proof}
  Project \(A\), \(\mathfrak m\), the map \(A\to T\), and the Noetherian,
  local, and domain fields from \cref{lem:jensen_selected_source_exists_source}.
\end{proof}

\begin{lemma}[Jensen selected completion equivalence]
  \label{lem:jensen_selected_ring_completion_equiv_source}
  \lean{Run202608192034.TODO.jensen_selected_ring_completion_equiv_source}
  \uses{lem:jensen_selected_source_exists_source,def:jensen_completion_witness}
  The selected source package carries the completion equivalence
  \(\widehat A\cong T\) together with the compatible map \(A\to T\).
\end{lemma}

\begin{proof}
  Project the `JensenCompletionWitness` field from the selected source package.
\end{proof}

\begin{lemma}[Jensen selected ring is a UFD]
  \label{lem:jensen_selected_ring_ufd_source}
  \lean{Run202608192034.TODO.jensen_selected_ring_ufd_source}
  \uses{lem:jensen_selected_source_exists_source,def:jensen_completion_witness}
  The selected source ring is a UFD.
\end{lemma}

\begin{proof}
  Project the UFD field stored in the selected source package's completion
  witness.
\end{proof}

\begin{lemma}[Jensen selected generic formal fiber]
  \label{lem:jensen_selected_ring_generic_fiber_source}
  \lean{Run202608192034.TODO.jensen_selected_ring_generic_fiber_source}
  \uses{lem:jensen_selected_source_exists_source}
  In the selected source ring, the zero prime of \(T\) contracts to zero, and
  every nonzero prime of \(T\) contracts nontrivially.
\end{lemma}

\begin{proof}
  Project the two contraction fields from the selected source package.
\end{proof}

\begin{lemma}[Jensen special case source package]
  \label{lem:jensen_special_case_source}
  \lean{Run202608192034.TODO.jensenSpecialCase_source}
  \uses{lem:jensen_selected_ring_exists_source,lem:jensen_selected_ring_completion_equiv_source,lem:jensen_selected_ring_ufd_source,lem:jensen_selected_ring_generic_fiber_source}
  There exists a two-dimensional Noetherian local UFD \(A\) with
  \(\widehat A\cong T\) such that the only prime ideal of \(T\) contracting
  to zero in \(A\) is \((0)\), and this selected source ring carries the
  completion witness used downstream.
\end{lemma}

\begin{proof}
  Assemble the selected ring, completion witness, UFD output, and contraction
  statement from the selected source package projections.
\end{proof}

\begin{lemma}[Jensen special case]
  \label{lem:jensen_special_case}
  \lean{Run202608192034.TODO.jensenSpecialCase}
  \uses{lem:jensen_special_case_source}
  The source package gives the counterexample-ring candidate used in the rest of
  the proof.
\end{lemma}

\begin{proof}
  Project the selected ring, map, contraction statement, and completion witness
  from \cref{lem:jensen_special_case_source}.
\end{proof}

\section{The bad quotient}

\begin{definition}[The counterexample candidate]
  \label{def:counterexample_ring}
  \lean{Run202608192034.TODO.counterexampleRing}
  \uses{lem:jensen_special_case}
  Fix a ring \(A\), a maximal ideal \(\mathfrak m_A\), and a completion
  isomorphism \(\widehat A\cong T\) supplied by
  \cref{lem:jensen_special_case}.
\end{definition}

\begin{proof}
  The preceding existence lemma supplies a witness; fix one such witness for
  the remainder of the construction.
\end{proof}

\begin{lemma}[Properties of the counterexample candidate]
  \label{lem:counterexample_ring_properties}
  \lean{Run202608192034.TODO.counterexampleRing_properties}
  \uses{def:counterexample_ring,lem:jensen_special_case}
  The ring \(A\) is a two-dimensional Noetherian local UFD and a domain, its
  completion is identified with \(T\), and every nonzero prime of \(T\) has
  nonzero contraction to \(A\).  Moreover, the completion map embeds \(A\)
  into \(T\).
\end{lemma}

\begin{proof}
  The structural properties and contraction assertion come from the chosen
  witness.  A UFD is a domain.  The kernel of the completion map is
  \(\bigcap_{n\geq1}\mathfrak m_A^n\), which is zero by the Krull intersection
  theorem for the Noetherian local domain \(A\).  Hence the completion
  isomorphism lets us regard \(A\) as a subring of \(T\).
\end{proof}

\begin{lemma}[The candidate is weakly quasi-complete]
  \label{lem:counterexample_ring_weak}
  \lean{Run202608192034.TODO.counterexampleRing_weaklyQuasiComplete}
  \uses{lem:counterexample_ring_properties,lem:farley_completion_criterion}
  The ring \(A\) is weakly quasi-complete.
\end{lemma}

\begin{proof}
  By \cref{lem:counterexample_ring_properties}, every nonzero prime of
  \(\widehat A\cong T\) contracts nontrivially to \(A\).  Apply
  \cref{lem:farley_completion_criterion}.
\end{proof}

\begin{definition}[Contraction of the distinguished prime]
  \label{def:contracted_prime}
  \lean{Run202608192034.TODO.contractedPrime}
  \uses{def:counterexample_ring,def:node_prime,lem:counterexample_ring_properties}
  Under the fixed embedding \(A\subseteq T\), define \(q=Q\cap A\).
\end{definition}

\begin{proof}
  This is a definition.
\end{proof}

\begin{lemma}[Height of the contracted prime]
  \label{lem:contracted_prime_height}
  \lean{Run202608192034.TODO.contractedPrime_nonzero_height_one}
  \uses{def:contracted_prime,lem:counterexample_ring_properties,lem:node_prime_height}
  The ideal \(q\) is a nonzero prime ideal of \(A\) of height one.
\end{lemma}

\begin{proof}
  Contraction preserves primality.  Since \(Q\neq(0)\) and the zero prime is
  the only prime of \(T\) contracting to zero, \(q\neq(0)\).  The completion
  map is faithfully flat, so going-down gives
  \(\operatorname{ht}q\leq\operatorname{ht}Q=1\).  A nonzero prime in the
  domain \(A\) has height at least one, hence equality holds.
\end{proof}

\begin{definition}[A prime generator]
  \label{def:prime_generator}
  \lean{Run202608192034.TODO.primeGenerator}
  \uses{lem:contracted_prime_height,lem:counterexample_ring_properties}
  Choose a prime element \(a\in A\) for which \(q=aA\).
\end{definition}

\begin{proof}
  Every height-one prime of the UFD \(A\) is principal and generated by a
  prime element, so \cref{lem:contracted_prime_height} supplies such \(a\).
\end{proof}

\begin{lemma}[Jensen's selected ring is a UFD]
  \label{lem:jensen_special_case_ufd_source}
  \lean{Run202608192034.TODO.jensenSpecialCase_isUFD_source}
  \uses{def:jensen_completion_witness}
  The ring \(A\) selected from Jensen's construction is a UFD.
\end{lemma}

\begin{proof}
  This is now a projection from the Jensen completion witness.  The witness
  records the UFD output of the saturated-union construction together with the
  completion equivalence and weak-completeness criterion.
\end{proof}

\begin{lemma}[Weak-completeness criterion for the selected ring]
  \label{lem:counterexample_weak_criterion_source}
  \lean{Run202608192034.TODO.counterexampleRing_weakCriterion_source}
  \uses{def:jensen_completion_witness}
  The selected ring \(A\) is weakly quasi-complete exactly when every nonzero
  prime of \(T\) contracts nontrivially to \(A\).
\end{lemma}

\begin{proof}
  This is now a projection from the Jensen completion witness, which records
  the Farley completion-prime criterion already transported along
  \(\widehat A\cong T\).
\end{proof}

\begin{lemma}[A nonzero ideal in a domain has positive height]
  \label{lem:nonzero_prime_height_ge_one_source}
  \lean{Run202608192034.TODO.nonzeroPrime_height_ge_one_source}
  If \(A\) is a domain and \(q\neq(0)\), then \(\operatorname{ht}q\geq 1\).
\end{lemma}

\begin{proof}
  The zero ideal is prime in a domain.  Since \((0)\subsetneq q\), strict
  monotonicity of height below a prime ideal gives
  \(\operatorname{ht}(0)<\operatorname{ht}q\), and
  \(\operatorname{ht}(0)=0\).
\end{proof}

\begin{lemma}[Going-down controls height under lies-over]
  \label{lem:lies_over_height_le_of_has_going_down_source}
  \lean{Run202608192034.TODO.liesOver_height_le_of_hasGoingDown_source}
  If \(T\) is a Noetherian \(A\)-algebra satisfying going-down and
  \(Q\subseteq T\) lies over \(q\subseteq A\), then
  \(\operatorname{ht}q\leq\operatorname{ht}Q\).
\end{lemma}

\begin{proof}
  Use Mathlib's height formula for going-down:
  \[
    \operatorname{ht}Q =
    \operatorname{ht}q +
    \operatorname{ht}\bigl(Q(T/qT)\bigr).
  \]
  Since \(x\leq x+y\) in \(\mathbb N_\infty\), the desired inequality follows.
\end{proof}

\begin{lemma}[Noetherian adic completions have going-down]
  \label{lem:adic_completion_has_going_down_of_is_noetherian}
  \lean{Run202608192034.TODO.adicCompletion_hasGoingDown_of_isNoetherian}
  If \(A\) is Noetherian, then the canonical algebra map
  \(A\to\widehat A^{\,\mathfrak m}\) has going-down.
\end{lemma}

\begin{proof}
  Mathlib proves that Noetherian adic completions are flat over the original
  ring.  Flat algebras satisfy going-down, so typeclass inference supplies the
  result.
\end{proof}

\begin{lemma}[Going-down transports across a completion equivalence]
  \label{lem:adic_completion_equiv_has_going_down_of_is_noetherian}
  \lean{Run202608192034.TODO.adicCompletion_equiv_hasGoingDown_of_isNoetherian}
  \uses{lem:adic_completion_has_going_down_of_is_noetherian}
  If \(A\) is Noetherian and \(\widehat A^{\,\mathfrak m}\cong T\), then the
  transported algebra structure \(A\to T\) has going-down.
\end{lemma}

\begin{proof}
  The canonical map \(A\to\widehat A^{\,\mathfrak m}\) is flat, and a bijective
  ring map \(\widehat A^{\,\mathfrak m}\to T\) is flat.  Flatness is stable
  under composition, hence the transported map \(A\to T\) is flat.  Apply
  `Algebra.HasGoingDown.of_flat`.
\end{proof}

\begin{lemma}[The completion map has going-down]
  \label{lem:completion_map_has_going_down_source}
  \lean{Run202608192034.TODO.completionMap_hasGoingDown_source}
  \uses{def:jensen_completion_witness,lem:adic_completion_equiv_has_going_down_of_is_noetherian}
  The selected completion map \(A\to T\) satisfies going-down.
\end{lemma}

\begin{proof}
  The Jensen completion witness identifies the selected map \(A\to T\) with
  the standard adic completion map transported across
  \(\widehat A\cong T\).  Apply the checked transport lemma for going-down
  across such a completion equivalence.
\end{proof}

\begin{lemma}[Going-down gives the contracted-prime height upper bound]
  \label{lem:contracted_prime_height_le_one_of_has_going_down_source}
  \lean{Run202608192034.TODO.contractedPrime_height_le_one_of_hasGoingDown_source}
  \uses{lem:lies_over_height_le_of_has_going_down_source,lem:node_prime_height}
  If the selected completion map has going-down, then
  \(\operatorname{ht}(Q\cap A)\leq1\).
\end{lemma}

\begin{proof}
  The equality \(q=Q\cap A\) says exactly that \(Q\) lies over \(q\).  Apply
  the going-down height comparison and then use
  \(\operatorname{ht}Q=1\).
\end{proof}

\begin{lemma}[Going-down gives the contracted-prime height upper bound]
  \label{lem:contracted_prime_height_le_one_source}
  \lean{Run202608192034.TODO.contractedPrime_height_le_one_source}
  \uses{lem:completion_map_has_going_down_source,lem:contracted_prime_height_le_one_of_has_going_down_source}
  For \(q=Q\cap A\), the completion map gives
  \(\operatorname{ht}q\leq 1\).
\end{lemma}

\begin{proof}
  Combine the source fact that the selected completion map has going-down with
  the checked height comparison for ideals satisfying lies-over.
\end{proof}

\begin{lemma}[The contracted prime has height one]
  \label{lem:contracted_prime_height_one_source}
  \lean{Run202608192034.TODO.contractedPrime_height_one_source}
  \uses{lem:nonzero_prime_height_ge_one_source,lem:contracted_prime_height_le_one_source}
  For \(q=Q\cap A\), the contracted prime \(q\) has height one.
\end{lemma}

\begin{proof}
  Combine the going-down upper bound \(\operatorname{ht}q\leq1\) with the
  lower bound \(\operatorname{ht}q\geq1\) coming from \(q\neq(0)\) in the
  domain \(A\).
\end{proof}

\begin{lemma}[Height-one primes in the selected UFD are principal]
  \label{lem:height_one_prime_principal_ufd_source}
  \lean{Run202608192034.TODO.heightOnePrime_principal_of_ufd_source}
  \uses{lem:jensen_special_case_ufd_source,lem:contracted_prime_height_one_source}
  A nonzero height-one prime ideal in the selected UFD is principal.
\end{lemma}

\begin{proof}
  Mathlib gives a prime element \(p\in q\) inside any nonzero prime ideal of a
  UFD.  Then \((p)\) is a nonzero prime ideal and \((p)\subseteq q\).  If the
  containment were strict, strict monotonicity of prime height would give
  \(\operatorname{ht}(p)<\operatorname{ht}q=1\), contradicting
  \(\operatorname{ht}(p)>0\).
\end{proof}

\begin{lemma}[Source theorem for the prime generator]
  \label{lem:prime_generator_source}
  \lean{Run202608192034.TODO.primeGenerator_source}
  \uses{def:prime_generator,lem:jensen_special_case,lem:counterexample_weak_criterion_source,lem:contracted_prime_height_one_source,lem:height_one_prime_principal_ufd_source,lem:jensen_special_case_ufd_source}
  The Jensen construction and the UFD height-one-prime argument supply the
  prime-generator package used later.
\end{lemma}

\begin{proof}
  Use the constructed ring \(A\), take \(q=Q\cap A\), and apply the
  height-one-prime-in-a-UFD argument to choose \(a\) with \(q=aA\).  The
  generic-fiber condition gives nonzeroness of the contraction.  The Jensen
  completion witness supplies both the transported Farley criterion and the
  going-down input used in the height calculation.
\end{proof}

\begin{lemma}[The chosen extended principal ideal is not prime]
  \label{lem:extended_principal_specific_not_prime}
  \lean{Run202608192034.TODO.extendedPrincipal_not_prime_of_generator_data}
  \uses{def:prime_generator,lem:counterexample_ring_properties,lem:node_prime_height,lem:node_prime_nonprincipal}
  For the specific generator \(a\) of the contracted prime \(q\), the extended
  ideal \(aT\) is not prime.
\end{lemma}

\begin{proof}
  The identity \(q=aA\) gives \(qT=aT\), while
  \(q=Q\cap A\) gives \(qT\subseteq Q\).  Since \(A\subseteq T\), the extended
  ideal is nonzero.  If \(aT\) were prime, the principal ideal theorem would
  give height at most one, and nonzeroness in the domain \(T\) gives height at
  least one.  A strict containment \(aT\subsetneq Q\) would then strictly
  increase prime height, contradicting that both have height one.  Hence
  \(aT=Q\), contradicting that \(Q\) is not principal.
\end{proof}

\begin{lemma}[The extended principal ideal is not prime]
  \label{lem:extended_principal_not_prime}
  \lean{Run202608192034.TODO.extendedPrincipal_not_prime}
  \uses{lem:extended_principal_specific_not_prime}
  The ideal \(aT\) is a proper nonprime ideal of \(T\).
\end{lemma}

\begin{proof}
  This is the existential form of
  \cref{lem:extended_principal_specific_not_prime}.
\end{proof}

\begin{definition}[The bad quotient]
  \label{def:bad_quotient}
  \lean{Run202608192034.TODO.badQuotient}
  \uses{def:counterexample_ring,def:prime_generator}
  Set \(B=A/aA\).
\end{definition}

\begin{proof}
  This is a definition.
\end{proof}

\begin{definition}[Bad quotient source data]
  \label{def:bad_quotient_source_data}
  \lean{Run202608192034.TODO.BadQuotientSourceData}
  \uses{def:bad_quotient,def:prime_generator}
  Package the source data for the bad quotient as one object:
  \(A\), \(\mathfrak m_A\), the map \(A\to T\), the contracted prime
  \(q=Q\cap A\), its chosen generator \(a\), and the structural properties of
  \(A\) and \(q\).
\end{definition}

\begin{proof}
  This is a data structure organizing the hypotheses already present in
  \cref{def:bad_quotient}, together with the Jensen generic-fiber data needed
  to recover the earlier contracted-prime and prime-generator nodes.
\end{proof}

\begin{definition}[Quotient criterion field]
  \label{def:bad_quotient_quasi_criterion_field}
  \lean{Run202608192034.TODO.BadQuotientSourceData.QuasiCriterion}
  \uses{def:bad_quotient_source_data,lem:quotient_weak_characterization}
  The source data carries the Anderson/Farley criterion saying that
  \(A\) is quasi-complete exactly when all quotients \(A/J\) are weakly
  quasi-complete.
\end{definition}

\begin{proof}
  This field records the criterion in the notation of the chosen data package.
\end{proof}

\begin{definition}[Dimension-one criterion field]
  \label{def:bad_quotient_dimension_criterion_field}
  \lean{Run202608192034.TODO.BadQuotientSourceData.DimensionCriterion}
  \uses{def:bad_quotient_source_data,lem:dimension_one_weak_criterion}
  The source data carries the dimension-one criterion identifying weak
  quasi-completeness of \(A/q\) with analytic irreducibility of the completed
  quotient \(T/aT\).
\end{definition}

\begin{proof}
  This field records the dimension-one criterion in the notation of the chosen
  data package.
\end{proof}

\begin{lemma}[The bad quotient is Noetherian]
  \label{lem:bad_quotient_source_data_quotient_is_noetherian}
  \lean{Run202608192034.TODO.BadQuotientSourceData.quotient_isNoetherian}
  \uses{def:bad_quotient_source_data}
  The quotient ring \(A/q\) is Noetherian.
\end{lemma}

\begin{proof}
  This follows from the stored Noetherian hypothesis on \(A\) and Mathlib's
  Noetherian quotient instance.
\end{proof}

\begin{lemma}[The bad quotient is a domain]
  \label{lem:bad_quotient_source_data_quotient_is_domain}
  \lean{Run202608192034.TODO.BadQuotientSourceData.quotient_isDomain}
  \uses{def:bad_quotient_source_data}
  The quotient ring \(A/q\) is a domain.
\end{lemma}

\begin{proof}
  Since \(q\) is prime, this is `Ideal.Quotient.isDomain_iff_prime`.
\end{proof}

\begin{lemma}[The bad quotient is nontrivial]
  \label{lem:bad_quotient_source_data_quotient_nontrivial}
  \lean{Run202608192034.TODO.BadQuotientSourceData.quotient_nontrivial}
  \uses{def:bad_quotient_source_data}
  The quotient ring \(A/q\) is nontrivial.
\end{lemma}

\begin{proof}
  A prime ideal is proper, so the quotient by \(q\) is nontrivial.
\end{proof}

\begin{lemma}[The quotient map is local]
  \label{lem:bad_quotient_source_data_quotient_mk_is_local_hom}
  \lean{Run202608192034.TODO.BadQuotientSourceData.quotient_mk_isLocalHom}
  \uses{lem:bad_quotient_source_data_quotient_nontrivial}
  The quotient map \(A\to A/q\) is a local homomorphism.
\end{lemma}

\begin{proof}
  Use `IsLocalHom.of_surjective` for the surjective quotient map.
\end{proof}

\begin{lemma}[The bad quotient is local]
  \label{lem:bad_quotient_source_data_quotient_is_local}
  \lean{Run202608192034.TODO.BadQuotientSourceData.quotient_isLocal}
  \uses{lem:bad_quotient_source_data_quotient_nontrivial,lem:bad_quotient_source_data_quotient_mk_is_local_hom}
  The quotient ring \(A/q\) is local.
\end{lemma}

\begin{proof}
  Apply `IsLocalRing.of_surjective` to the quotient map.
\end{proof}

\begin{lemma}[The bad quotient is a Noetherian local domain]
  \label{lem:bad_quotient_source_data_quotient_noetherian_local_domain}
  \lean{Run202608192034.TODO.BadQuotientSourceData.quotient_noetherian_local_domain}
  \uses{lem:bad_quotient_source_data_quotient_is_noetherian,lem:bad_quotient_source_data_quotient_is_local,lem:bad_quotient_source_data_quotient_is_domain}
  The quotient ring \(A/q\) is a Noetherian local domain.
\end{lemma}

\begin{proof}
  Combine the separately proved Noetherian, local, and domain facts.
\end{proof}

\begin{lemma}[Recover the Jensen completion witness from bad quotient data]
  \label{lem:bad_quotient_source_data_jensen_completion_witness}
  \lean{Run202608192034.TODO.BadQuotientSourceData.jensenCompletionWitness}
  \uses{def:bad_quotient_source_data,def:jensen_completion_witness}
  The bad quotient source package determines the Jensen completion witness
  \(\widehat A\cong T\) for the chosen map \(A\to T\).
\end{lemma}

\begin{proof}
  Project the stored Jensen completion witness from the bad quotient source
  data.
\end{proof}

\begin{lemma}[The Jensen completion is Noetherian]
  \label{lem:adic_completion_is_noetherian_ring_from_jensen}
  \lean{Run202608192034.TODO.adicCompletion_isNoetherianRing_from_jensen}
  \uses{def:jensen_completion_witness,lem:node_complete_cm_dim,lem:is_noetherian_ring_of_ring_equiv_source}
  If Jensen's construction identifies \(\widehat A\) with the chosen node
  ring \(T\), then \(\widehat A\) is Noetherian.
\end{lemma}

\begin{proof}
  The complete node ring \(T\) is Noetherian by the node-ring package.
  Transport this Noetherianity backward across the inverse of the Jensen
  completion equivalence \(\widehat A\cong T\).
\end{proof}

\begin{lemma}[The source ring is two-dimensional]
  \label{lem:bad_quotient_source_data_source_ring_krull_dim_eq_two}
  \lean{Run202608192034.TODO.BadQuotientSourceData.source_ringKrullDim_eq_two}
  \uses{lem:bad_quotient_source_data_jensen_completion_witness}
  The Jensen source ring \(A\) has Krull dimension two.
\end{lemma}

\begin{proof}
  This is the dimension field of the Jensen completion witness: the selected
  \(A\) has completion \(T\), and \(T\) is the two-dimensional node ring.
\end{proof}

\begin{lemma}[The source completion is Noetherian]
  \label{lem:bad_quotient_source_data_adic_completion_is_noetherian_ring}
  \lean{Run202608192034.TODO.BadQuotientSourceData.adicCompletion_isNoetherianRing}
  \uses{lem:bad_quotient_source_data_jensen_completion_witness,lem:adic_completion_is_noetherian_ring_from_jensen}
  For the bad quotient source data, the completion
  \(\widehat A^{\,\mathfrak m}\) is Noetherian.
\end{lemma}

\begin{proof}
  Apply the Jensen completion Noetherianity lemma to the completion witness
  recovered from the source data.
\end{proof}

\begin{lemma}[The contracted prime has height one in the source ring]
  \label{lem:bad_quotient_source_data_q_height_one_source}
  \lean{Run202608192034.TODO.BadQuotientSourceData.q_height_one_source}
  \uses{lem:bad_quotient_source_data_jensen_completion_witness,lem:contracted_prime_height_one_source}
  The contracted prime \(q=Q\cap A\) has height one.
\end{lemma}

\begin{proof}
  Reuse the already proved contracted-prime height calculation, using the
  Jensen completion witness recovered from the source data.
\end{proof}

\begin{lemma}[The chosen generator lies in the contracted prime]
  \label{lem:bad_quotient_source_data_generator_mem_q}
  \lean{Run202608192034.TODO.BadQuotientSourceData.generator_mem_q}
  \uses{def:bad_quotient_source_data}
  The chosen generator \(a\) lies in \(q\).
\end{lemma}

\begin{proof}
  Rewrite \(q=(a)\) and use membership of a generator in its principal ideal.
\end{proof}

\begin{lemma}[The chosen generator is nonzero]
  \label{lem:bad_quotient_source_data_generator_ne_zero}
  \lean{Run202608192034.TODO.BadQuotientSourceData.generator_ne_zero}
  \uses{def:bad_quotient_source_data}
  The chosen generator \(a\) is nonzero.
\end{lemma}

\begin{proof}
  If \(a=0\), then \(q=(a)=0\), contradicting the stored nonzero property of
  \(q\).
\end{proof}

\begin{lemma}[The chosen generator is a non-zero-divisor]
  \label{lem:bad_quotient_source_data_generator_mem_non_zero_divisors}
  \lean{Run202608192034.TODO.BadQuotientSourceData.generator_mem_nonZeroDivisors}
  \uses{lem:bad_quotient_source_data_generator_ne_zero}
  The chosen generator \(a\) belongs to the non-zero-divisors of \(A\).
\end{lemma}

\begin{proof}
  Since \(A\) is a domain, every nonzero element is a non-zero-divisor.
\end{proof}

\begin{lemma}[The chosen generator lies in the maximal ideal]
  \label{lem:bad_quotient_source_data_generator_mem_maximal_ideal}
  \lean{Run202608192034.TODO.BadQuotientSourceData.generator_mem_maximalIdeal}
  \uses{lem:bad_quotient_source_data_generator_mem_q}
  The chosen generator \(a\) lies in the maximal ideal of the local ring \(A\).
\end{lemma}

\begin{proof}
  The proper prime ideal \(q\) is contained in the maximal ideal of the local
  ring \(A\), and \(a\in q\).
\end{proof}

\begin{lemma}[Quotient by the generator lowers dimension by one]
  \label{lem:bad_quotient_source_data_quotient_span_generator_dimension_add_one}
  \lean{Run202608192034.TODO.BadQuotientSourceData.quotient_span_generator_dimension_add_one}
  \uses{lem:bad_quotient_source_data_generator_mem_non_zero_divisors,lem:bad_quotient_source_data_generator_mem_maximal_ideal}
  The dimension of \(A/(a)\) satisfies
  \[
    \dim(A/(a))+1=\dim A.
  \]
\end{lemma}

\begin{proof}
  Apply Mathlib's dimension-drop theorem for quotienting a Noetherian local
  ring by a non-zero-divisor lying in the maximal ideal.
\end{proof}

\begin{lemma}[The quotient dimension adds to two]
  \label{lem:bad_quotient_source_data_quotient_dimension_add_one_eq_two_source}
  \lean{Run202608192034.TODO.BadQuotientSourceData.quotient_dimension_add_one_eq_two_source}
  \uses{lem:bad_quotient_source_data_quotient_span_generator_dimension_add_one,lem:bad_quotient_source_data_source_ring_krull_dim_eq_two}
  The dimension of \(A/q\) satisfies
  \[
    \dim(A/q)+1=2.
  \]
\end{lemma}

\begin{proof}
  Rewrite \(q=(a)\) in the previous dimension-drop lemma, then use
  \(\dim A=2\).
\end{proof}

\begin{lemma}[WithBot ENat cancellation at two]
  \label{lem:with_bot_enat_eq_one_of_add_one_eq_two}
  \lean{Run202608192034.TODO.withBotENat_eq_one_of_add_one_eq_two}
  If \(x+1=2\) in \(\bot\)-extended \(\mathbb N_\infty\), then \(x=1\).
\end{lemma}

\begin{proof}
  Rewrite \(2\) as \(1+1\) and use the order-theoretic cancellation lemma for
  adding one on both sides.
\end{proof}

\begin{lemma}[The bad quotient has dimension one]
  \label{lem:bad_quotient_source_data_quotient_ring_krull_dim_eq_one_source}
  \lean{Run202608192034.TODO.BadQuotientSourceData.quotient_ringKrullDim_eq_one_source}
  \uses{lem:bad_quotient_source_data_quotient_dimension_add_one_eq_two_source,lem:with_bot_enat_eq_one_of_add_one_eq_two}
  The quotient \(A/q\) has Krull dimension one.
\end{lemma}

\begin{proof}
  Apply the extended-natural-number cancellation lemma to
  \(\dim(A/q)+1=2\).
\end{proof}

\begin{definition}[Dimension-one facts for the bad quotient]
  \label{def:quotient_dimension_one_facts}
  \lean{Run202608192034.TODO.QuotientDimensionOneFacts}
  \uses{lem:bad_quotient_source_data_quotient_noetherian_local_domain,lem:bad_quotient_source_data_quotient_ring_krull_dim_eq_one_source}
  Record the exact hypotheses needed before applying the one-dimensional
  quotient criterion: \(A/q\) is Noetherian, local, a domain, and has Krull
  dimension one.
\end{definition}

\begin{proof}
  This is a structure definition isolating the local algebra facts used by the
  one-dimensional criterion.
\end{proof}

\begin{lemma}[Dimension-one facts from the dimension calculation]
  \label{lem:bad_quotient_source_data_quotient_dimension_one_facts_of_dim}
  \lean{Run202608192034.TODO.BadQuotientSourceData.quotientDimensionOneFacts_of_dim}
  \uses{def:quotient_dimension_one_facts,lem:bad_quotient_source_data_quotient_noetherian_local_domain,lem:bad_quotient_source_data_quotient_ring_krull_dim_eq_one_source}
  Once the source proof gives \(\dim(A/q)=1\), all one-dimensional quotient
  facts are available.
\end{lemma}

\begin{proof}
  Fill the structure using the checked Noetherian, local, and domain facts,
  and the supplied dimension equality.
\end{proof}

\begin{definition}[Jensen completion witness]
  \label{def:jensen_completion_witness}
  \lean{Run202608192034.TODO.JensenCompletionWitness}
  \uses{lem:farley_completion_criterion}
  Record the stronger Jensen output needed later: an abstract completion
  equivalence \(\widehat A^{\,\mathfrak m}\cong T\), compatibility with the
  embedded map \(A\to T\), the UFD property of the selected source ring, the
  fact that the chosen adic ideal is the maximal ideal of the local ring \(A\),
  and the weak quasi-completeness criterion expressed through this completion.
\end{definition}

\begin{proof}
  This is a source data structure isolating the part of Jensen's construction
  that identifies the completion of \(A\) with the chosen complete local ring
  \(T\).  In Lean the completion is the standard Mathlib object
  `AdicCompletion \(\mathfrak m\) A`; the maximal-ideal equality records the
  paper's implicit notation \((A,\mathfrak m)\).
\end{proof}

\begin{definition}[Quotient completion witness]
  \label{def:quotient_completion_witness}
  \lean{Run202608192034.TODO.QuotientCompletionWitness}
  \uses{def:bad_quotient_source_data,def:jensen_completion_witness,lem:dimension_one_weak_criterion}
  Record the completion bridge for the quotient \(A/q=A/aA\): its completion
  target is equivalent to \(T/aT\), and the dimension-one analytic
  irreducibility criterion applies to that completion target.
\end{definition}

\begin{proof}
  This is the source interface for the theorem that completion commutes with
  quotient by the principal ideal \(aA\), after transporting
  \(\widehat A\cong T\).
\end{proof}

\begin{definition}[The quotient map as a linear map]
  \label{def:quotient_linear_map}
  \lean{Run202608192034.TODO.quotientLinearMap}
  \uses{def:bad_quotient_source_data}
  View the canonical quotient map \(A\to A/q\) as an \(A\)-linear map, so that
  Mathlib's module-level adic-completion functoriality can be applied.
\end{definition}

\begin{proof}
  This is the linear map underlying the algebra homomorphism
  `Ideal.Quotient.mkₐ`.
\end{proof}

\begin{lemma}[The quotient map is surjective]
  \label{lem:quotient_mk_linear_surjective}
  \lean{Run202608192034.TODO.quotient_mk_linear_surjective}
  \uses{def:quotient_linear_map}
  The linear version of the quotient map \(A\to A/q\) is surjective.
\end{lemma}

\begin{proof}
  This is exactly `Ideal.Quotient.mk_surjective`, after unfolding the linear
  map.
\end{proof}

\begin{lemma}[Completion preserves the quotient-map surjection]
  \label{lem:quotient_completion_map_surjective}
  \lean{Run202608192034.TODO.quotient_completion_map_surjective}
  \uses{def:quotient_linear_map,lem:quotient_mk_linear_surjective}
  The completed map induced by \(A\to A/q\) is surjective.
\end{lemma}

\begin{proof}
  Apply Mathlib's `AdicCompletion.map_surjective` to the surjective quotient
  linear map.
\end{proof}

\begin{lemma}[The two quotient filtrations agree]
  \label{lem:quotient_module_filtration_eq_restrict_scalars}
  \lean{Run202608192034.TODO.quotient_module_filtration_eq_restrictScalars}
  \uses{def:bad_quotient_source_data}
  The \(A\)-module filtration on \(A/q\) induced by \(\mathfrak m^n\) agrees
  with the restriction of scalars of the ideal
  \((\mathfrak m(A/q))^n\).
\end{lemma}

\begin{proof}
  Rewrite the scalar action of an ideal on the top submodule as ideal image,
  identify the algebra map \(A\to A/q\) with the quotient map, and use
  `Ideal.map_pow`.
\end{proof}

\begin{definition}[Finite-level module-to-ring quotient equivalence]
  \label{def:quotient_module_level_equiv}
  \lean{Run202608192034.TODO.quotientModuleLevelEquiv}
  \uses{lem:quotient_module_filtration_eq_restrict_scalars}
  At each finite level, quotienting \(A/q\) by the \(A\)-module filtration
  \(\mathfrak m^n(A/q)\) is linearly equivalent to quotienting by the same
  ideal viewed after restriction of scalars from \(A/q\).
\end{definition}

\begin{proof}
  Apply `Submodule.Quotient.equiv` to the equality of submodules from
  \cref{lem:quotient_module_filtration_eq_restrict_scalars}.
\end{proof}

\begin{lemma}[The finite-level equivalence on representatives]
  \label{lem:quotient_module_level_equiv_mk}
  \lean{Run202608192034.TODO.quotientModuleLevelEquiv_mk}
  \uses{def:quotient_module_level_equiv}
  The finite-level equivalence sends the class of \(x\in A/q\) to the class of
  the same representative \(x\).
\end{lemma}

\begin{proof}
  This is definitional for the quotient equivalence built from the identity
  linear equivalence.
\end{proof}

\begin{definition}[Finite-level comparison to the ring quotient]
  \label{def:quotient_module_level_to_ring_level_equiv}
  \lean{Run202608192034.TODO.quotientModuleLevelToRingLevelEquiv}
  \uses{def:quotient_module_level_equiv}
  This is the same finite-level comparison, with the target written as the
  ordinary ideal quotient over the ring \(A/q\).
\end{definition}

\begin{proof}
  It is the previous equivalence after Lean unfolds ideals as submodules of
  their underlying ring.
\end{proof}

\begin{lemma}[The module filtration agrees with the adic filtration]
  \label{lem:quotient_module_filtration_eq_adic_restrict_scalars}
  \lean{Run202608192034.TODO.quotient_module_filtration_eq_adic_restrictScalars}
  \uses{lem:quotient_module_filtration_eq_restrict_scalars}
  The \(A\)-module filtration on \(A/q\) agrees with the restricted scalar
  version of the \(A/q\)-adic filtration.
\end{lemma}

\begin{proof}
  Combine the already established filtration equality with the fact that an
  ideal acting on the top submodule of its own ring gives the ideal itself.
\end{proof}

\begin{definition}[Finite-level comparison to adic levels]
  \label{def:quotient_module_level_to_adic_level_equiv}
  \lean{Run202608192034.TODO.quotientModuleLevelToAdicLevelEquiv}
  \uses{lem:quotient_module_filtration_eq_adic_restrict_scalars}
  At finite level \(n\), the \(A\)-module completion quotient of \(A/q\) is
  linearly equivalent to the \(n\)-th quotient used in the intrinsic
  \(A/q\)-adic completion.
\end{definition}

\begin{proof}
  Apply `Submodule.Quotient.equiv` to the equality of restricted scalar
  submodules.
\end{proof}

\begin{lemma}[Compatibility with the ring-level factor map]
  \label{lem:quotient_module_level_to_adic_level_equiv_eq_factor_ring_level}
  \lean{Run202608192034.TODO.quotientModuleLevelToAdicLevelEquiv_eq_factor_ringLevel}
  \uses{def:quotient_module_level_to_adic_level_equiv,def:quotient_module_level_to_ring_level_equiv}
  The finite-level comparison to adic levels agrees with first passing through
  the ring-level quotient and then applying the canonical factor map.
\end{lemma}

\begin{proof}
  Prove the statement by quotient induction on representatives; both sides
  reduce to the same quotient class.
\end{proof}

\begin{lemma}[The completed contracted prime is principal]
  \label{lem:completed_q_eq_span_a}
  \lean{Run202608192034.TODO.completed_q_eq_span_a}
  \uses{def:bad_quotient_source_data}
  Since \(q=aA\), its extension to \(\widehat A\) is generated by the image of
  \(a\).
\end{lemma}

\begin{proof}
  Rewrite by the stored equality \(q=(a)\), then use `Ideal.map_span` and the
  singleton image simplification.
\end{proof}

\begin{lemma}[Transporting the completed principal ideal to \(T\)]
  \label{lem:completion_equiv_maps_completed_q}
  \lean{Run202608192034.TODO.completionEquiv_maps_completed_q}
  \uses{def:jensen_completion_witness,lem:completed_q_eq_span_a}
  The Jensen completion equivalence sends \(q\widehat A\) to the principal
  ideal \(aT\subseteq T\).
\end{lemma}

\begin{proof}
  Use `completed_q_eq_span_a` and the compatibility equation recorded in the
  Jensen completion witness, namely that the selected embedding \(A\to T\) is
  the standard completion map followed by \(\widehat A\cong T\).
\end{proof}

\begin{definition}[The quotient target transported to \(T/aT\)]
  \label{def:completed_quotient_target_equiv_node}
  \lean{Run202608192034.TODO.completedQuotientTargetEquivNode}
  \uses{def:jensen_completion_witness,lem:completion_equiv_maps_completed_q}
  The completion equivalence \(\widehat A\cong T\) induces a quotient
  equivalence \(\widehat A/q\widehat A\cong T/aT\).
\end{definition}

\begin{proof}
  Apply Mathlib's quotient equivalence attached to a ring equivalence and an
  equality of mapped ideals.
\end{proof}

\begin{definition}[The actual quotient adic completion]
  \label{def:quotient_adic_completion}
  \lean{Run202608192034.TODO.quotientAdicCompletion}
  \uses{def:bad_quotient_source_data}
  The completion of \(B=A/q\) is represented as the adic completion over the
  quotient ring \(A/q\) itself, with respect to the ideal induced by
  \(\mathfrak m\).
\end{definition}

\begin{proof}
  The Lean definition fixes the hidden base-ring and module parameters of
  `AdicCompletion` explicitly, avoiding confusion with the completion of
  \(A/q\) merely as an \(A\)-module.
\end{proof}

\begin{lemma}[The quotient adic completion is a ring]
  \label{lem:quotient_adic_completion_comm_ring}
  \lean{Run202608192034.TODO.quotientAdicCompletion.instCommRing}
  \uses{def:quotient_adic_completion}
  The quotient completion \( \widehat{A/q}\) has its natural commutative ring
  structure.
\end{lemma}

\begin{proof}
  This is Mathlib's `AdicCompletion.instCommRing`, applied with \(A/q\) as the
  base ring.
\end{proof}

\begin{definition}[From module completion to quotient adic completion]
  \label{def:quotient_module_completion_to_quotient_adic_completion}
  \lean{Run202608192034.TODO.quotientModuleCompletionToQuotientAdicCompletion}
  \uses{def:quotient_adic_completion,def:quotient_module_level_to_adic_level_equiv}
  The finite-level comparison maps assemble into a map from the completion of
  \(A/q\) as an \(A\)-module to the intrinsic adic completion of \(A/q\).
\end{definition}

\begin{proof}
  Send a compatible family levelwise through
  \cref{def:quotient_module_level_to_adic_level_equiv}; compatibility follows
  by reducing the transition-map identity to representatives.
\end{proof}

\begin{definition}[From quotient adic completion to module completion]
  \label{def:quotient_adic_completion_to_quotient_module_completion}
  \lean{Run202608192034.TODO.quotientAdicCompletionToQuotientModuleCompletion}
  \uses{def:quotient_adic_completion,def:quotient_module_level_to_adic_level_equiv}
  The inverse levelwise comparison maps assemble into a map from the intrinsic
  quotient adic completion back to the \(A\)-module completion.
\end{definition}

\begin{proof}
  Send a compatible family levelwise through the inverse finite-level
  equivalence; the transition-map check again reduces to representatives.
\end{proof}

\begin{lemma}[The comparison has a left inverse]
  \label{lem:quotient_module_completion_to_quotient_adic_completion_left_inverse}
  \lean{Run202608192034.TODO.quotientModuleCompletionToQuotientAdicCompletion_leftInverse}
  \uses{def:quotient_module_completion_to_quotient_adic_completion,def:quotient_adic_completion_to_quotient_module_completion}
  Going from module completion to quotient adic completion and back is the
  identity.
\end{lemma}

\begin{proof}
  Use extensionality of adic completions and the left inverse law for each
  finite-level linear equivalence.
\end{proof}

\begin{lemma}[The comparison has a right inverse]
  \label{lem:quotient_module_completion_to_quotient_adic_completion_right_inverse}
  \lean{Run202608192034.TODO.quotientModuleCompletionToQuotientAdicCompletion_rightInverse}
  \uses{def:quotient_module_completion_to_quotient_adic_completion,def:quotient_adic_completion_to_quotient_module_completion}
  Going from quotient adic completion to module completion and back is the
  identity.
\end{lemma}

\begin{proof}
  Use extensionality of adic completions and the right inverse law for each
  finite-level linear equivalence.
\end{proof}

\begin{definition}[The comparison equivalence of the two completions]
  \label{def:quotient_module_completion_equiv_quotient_adic_completion}
  \lean{Run202608192034.TODO.quotientModuleCompletionEquivQuotientAdicCompletion}
  \uses{lem:quotient_module_completion_to_quotient_adic_completion_left_inverse,lem:quotient_module_completion_to_quotient_adic_completion_right_inverse}
  The two Lean objects
  `AdicCompletion \(\mathfrak m\) (A/q)` and
  `quotientAdicCompletion` are equivalent.
\end{definition}

\begin{proof}
  Package the two comparison maps and their inverse laws as an equivalence.
\end{proof}

\begin{lemma}[The comparison map is surjective]
  \label{lem:quotient_module_completion_to_quotient_adic_completion_surjective}
  \lean{Run202608192034.TODO.quotientModuleCompletionToQuotientAdicCompletion_surjective}
  \uses{lem:quotient_module_completion_to_quotient_adic_completion_right_inverse}
  The map from the module completion of \(A/q\) to the intrinsic quotient adic
  completion is surjective.
\end{lemma}

\begin{proof}
  Surjectivity follows from the explicit right inverse.
\end{proof}

\begin{lemma}[Applying the dimension-one criterion source]
  \label{lem:quotient_dimension_one_facts_dimension_one_criterion_of_source}
  \lean{Run202608192034.TODO.QuotientDimensionOneFacts.dimensionOneCriterion_of_source}
  \uses{def:quotient_dimension_one_facts,def:quotient_adic_completion,lem:dimension_one_weak_criterion}
  Given the recorded one-dimensional quotient facts and the cited analytic
  criterion, weak quasi-completeness of \(A/q\) is equivalent to analytic
  irreducibility of its actual adic completion.
\end{lemma}

\begin{proof}
  Apply the project-level wrapper for the Anderson/Farley dimension-one
  criterion to the supplied source criterion.
\end{proof}

\begin{lemma}[The quotient criterion from a dimension calculation]
  \label{lem:quotient_dimension_one_criterion_on_quotient_completion_of_source}
  \lean{Run202608192034.TODO.quotient_dimensionOneCriterionOnQuotientCompletion_of_source}
  \uses{lem:bad_quotient_source_data_quotient_dimension_one_facts_of_dim,lem:quotient_dimension_one_facts_dimension_one_criterion_of_source}
  If the source proof supplies \(\dim(A/q)=1\) and the one-dimensional analytic
  criterion, then the required quotient-completion criterion follows.
\end{lemma}

\begin{proof}
  Build the dimension-one fact package from the dimension equality, then apply
  the criterion wrapper.
\end{proof}

\begin{lemma}[The chosen adic ideal is the maximal ideal]
  \label{lem:bad_quotient_source_data_adic_ideal_eq_maximal_ideal_source}
  \lean{Run202608192034.TODO.BadQuotientSourceData.adicIdeal_eq_maximalIdeal_source}
  \uses{def:bad_quotient_source_data}
  The ideal \(\mathfrak m\) carried by the bad-quotient source data is the
  maximal ideal of the local ring \(A\).
\end{lemma}

\begin{proof}
  This is now a checked projection from the Jensen completion witness.  The
  source proof of that witness records the paper's implicit notation
  \((A,\mathfrak m)\).
\end{proof}

\begin{lemma}[Project-specific Hausdorff quotient of the completion]
  \label{lem:bad_quotient_source_data_adic_completion_quotient_hausdorff_source}
  \lean{Run202608192034.TODO.BadQuotientSourceData.adicCompletion_quotient_hausdorff_source}
  \uses{lem:bad_quotient_source_data_adic_completion_is_noetherian_ring,lem:adic_completion_quotient_hausdorff_of_noetherian_completion}
  For the Jensen source data, every quotient \(\widehat A/P\) is Hausdorff
  for the \(\mathfrak m\widehat A\)-adic topology.
\end{lemma}

\begin{proof}
  Use the Jensen-specific Noetherianity of \(\widehat A\), then apply the
  Basic-level quotient-Hausdorff theorem conditional on Noetherianity of the
  completion.
\end{proof}

\begin{lemma}[Project-specific completion-neighborhood intersection]
  \label{lem:bad_quotient_source_data_adic_completion_neighborhood_iinf_le_prime_source}
  \lean{Run202608192034.TODO.BadQuotientSourceData.adicCompletion_neighborhood_iInf_le_prime_source}
  \uses{lem:bad_quotient_source_data_adic_completion_quotient_hausdorff_source,lem:adic_completion_neighborhood_iinf_le_prime_of_hausdorff}
  For the Jensen source data,
  \(\bigcap_n(P+\mathfrak m^n\widehat A)\subseteq P\).
\end{lemma}

\begin{proof}
  Install the project-specific quotient-Hausdorff instance and apply the
  checked Basic-level Hausdorff-to-intersection lemma.
\end{proof}

\begin{lemma}[Project-specific contracted-chain intersection containment]
  \label{lem:bad_quotient_source_data_bad_prime_contraction_chain_iinf_le_comap_source}
  \lean{Run202608192034.TODO.BadQuotientSourceData.badPrimeContractionChain_iInf_le_comap_source}
  \uses{lem:mem_bad_prime_contraction_chain,lem:bad_quotient_source_data_adic_completion_neighborhood_iinf_le_prime_source}
  For the Jensen source data,
  \(\bigcap_n I_n\subseteq P\cap A\) for Farley's contraction chain.
\end{lemma}

\begin{proof}
  Send an element of \(\bigcap I_n\) into \(\widehat A\), use the membership
  characterization of \(I_n\), and then apply the project-specific
  completion-neighborhood intersection containment.
\end{proof}

\begin{lemma}[Project-specific contracted-chain intersection equality]
  \label{lem:bad_quotient_source_data_bad_prime_contraction_chain_iinf_eq_comap_source}
  \lean{Run202608192034.TODO.BadQuotientSourceData.badPrimeContractionChain_iInf_eq_comap_source}
  \uses{lem:bad_quotient_source_data_bad_prime_contraction_chain_iinf_le_comap_source,lem:comap_le_bad_prime_contraction_chain_iinf}
  For the Jensen source data, Farley's contraction chain satisfies
  \(\bigcap_n I_n=P\cap A\).
\end{lemma}

\begin{proof}
  Combine the project-specific containment with the previously checked
  general reverse containment.
\end{proof}

\begin{lemma}[Project-specific contracted-chain intersection is zero]
  \label{lem:bad_quotient_source_data_bad_prime_contraction_chain_iinf_eq_bot_source}
  \lean{Run202608192034.TODO.BadQuotientSourceData.badPrimeContractionChain_iInf_eq_bot_source}
  \uses{lem:bad_quotient_source_data_bad_prime_contraction_chain_iinf_eq_comap_source}
  If a completion prime \(P\) has zero contraction, then Farley's contraction
  chain has zero intersection.
\end{lemma}

\begin{proof}
  Rewrite the project-specific chain intersection as \(P\cap A\), then use the
  zero-contraction hypothesis.
\end{proof}

\begin{lemma}[Project-specific bad completion prime gives a bad chain]
  \label{lem:bad_quotient_source_data_adic_completion_bad_prime_to_bad_chain_source}
  \lean{Run202608192034.TODO.BadQuotientSourceData.adicCompletion_badPrime_to_badChain_source}
  \uses{lem:bad_quotient_source_data_bad_prime_contraction_chain_iinf_eq_bot_source,lem:bad_prime_contraction_chain_antitone,lem:bad_prime_contraction_chain_avoids_fixed_power,lem:bad_quotient_source_data_adic_ideal_eq_maximal_ideal_source}
  For the selected Jensen source data, a bad completion prime produces
  Farley's forbidden descending chain.
\end{lemma}

\begin{proof}
  Use the same explicit chain \(I_n=(P+\mathfrak m^n\widehat A)\cap A\).
  The antitonicity and fixed-power avoidance are the checked Basic lemmas; the
  zero-intersection part now uses the project-specific Hausdorff quotient
  route instead of the remaining generic source theorem.
\end{proof}

\begin{lemma}[Project-specific weak completeness gives prime contraction]
  \label{lem:bad_quotient_source_data_weakly_quasi_complete_to_adic_completion_prime_contraction_source}
  \lean{Run202608192034.TODO.BadQuotientSourceData.weaklyQuasiComplete_to_adicCompletionPrimeContraction_source}
  \uses{lem:bad_quotient_source_data_adic_completion_bad_prime_to_bad_chain_source,lem:not_weakly_quasi_complete_iff_bad_chain,lem:not_adic_completion_prime_contraction_condition_iff_bad_prime}
  For the selected Jensen source data, weak quasi-completeness implies the
  adic-completion prime-contraction condition.
\end{lemma}

\begin{proof}
  Argue by contradiction.  Failure of prime contraction gives a bad completion
  prime; the project-specific previous lemma gives a bad descending chain; the
  checked negation lemma says this contradicts the stored weak
  quasi-completeness of \(A\).
\end{proof}

\begin{lemma}[Maximal ideals map to maximal ideals under the quotient map]
  \label{lem:bad_quotient_source_data_quotient_adic_ideal_eq_maximal_ideal_of_source}
  \lean{Run202608192034.TODO.BadQuotientSourceData.quotient_adicIdeal_eq_maximalIdeal_of_source}
  \uses{lem:bad_quotient_source_data_adic_ideal_eq_maximal_ideal_source,lem:bad_quotient_source_data_quotient_is_local}
  If \(\mathfrak m\) is the maximal ideal of \(A\), then its image in \(A/q\)
  is the maximal ideal of the quotient local ring.
\end{lemma}

\begin{proof}
  The quotient map \(A\to A/q\) is a surjective local homomorphism.  Apply
  Mathlib's `IsLocalRing.map_maximalIdeal_of_surjective`, then rewrite the
  source maximal ideal by the supplied equality.
\end{proof}

\begin{lemma}[The quotient adic ideal is the quotient maximal ideal]
  \label{lem:bad_quotient_source_data_quotient_adic_ideal_eq_maximal_ideal_source}
  \lean{Run202608192034.TODO.BadQuotientSourceData.quotient_adicIdeal_eq_maximalIdeal_source}
  \uses{lem:bad_quotient_source_data_adic_ideal_eq_maximal_ideal_source,lem:bad_quotient_source_data_quotient_adic_ideal_eq_maximal_ideal_of_source}
  The ideal \(\mathfrak m(A/q)\) used for the adic completion of \(A/q\) is
  the maximal ideal of the quotient local ring.
\end{lemma}

\begin{proof}
  Apply the checked quotient transport lemma to the recorded source equality
  \(\mathfrak m=\operatorname{maximalIdeal}(A)\).
\end{proof}

\begin{lemma}[Powers of the quotient maximal ideal]
  \label{lem:quotient_map_sup_power}
  \lean{Run202608192034.TODO.quotient_map_sup_power}
  \uses{def:bad_quotient_source_data}
  For every \(n\), the image of \(q+\mathfrak m^n\) in \(A/q\) is the
  \(n\)-th power of the image of \(\mathfrak m\).
\end{lemma}

\begin{proof}
  Use `Ideal.map_sup`, the fact that \(q\) maps to zero in \(A/q\), and
  `Ideal.map_pow`.
\end{proof}

\begin{definition}[Finite-level quotient equivalence for \(q+\mathfrak m^n\)]
  \label{def:quotient_power_quotient_equiv}
  \lean{Run202608192034.TODO.quotientPowerQuotientEquiv}
  \uses{lem:quotient_sup_quotient_equiv,lem:quotient_map_sup_power}
  For every \(n\), there is a natural equivalence
  \[
    A/(q+\mathfrak m^n)\cong
      (A/q)/(\mathfrak m(A/q))^n.
  \]
\end{definition}

\begin{proof}
  Apply the general quotient-by-sum equivalence to \(I=q\) and
  \(J=\mathfrak m^n\), then rewrite the image of \(\mathfrak m^n\) as the
  \(n\)-th power of the image of \(\mathfrak m\).
\end{proof}

\begin{lemma}[Finite-level equivalence on representatives]
  \label{lem:quotient_power_quotient_equiv_mk}
  \lean{Run202608192034.TODO.quotientPowerQuotientEquiv_mk}
  \uses{def:quotient_power_quotient_equiv}
  The finite-level equivalence sends the class of \(x\in A\) modulo
  \(q+\mathfrak m^n\) to the class of the class of \(x\) in \(A/q\).
\end{lemma}

\begin{proof}
  This is the representative formula for the general quotient-by-sum
  equivalence, followed by the power rewrite.
\end{proof}

\begin{definition}[Finite components of the map to the quotient completion]
  \label{def:completed_ring_to_quotient_finite}
  \lean{Run202608192034.TODO.completedRingToQuotientFinite}
  \uses{def:quotient_power_quotient_equiv}
  For every \(n\), the completed ring \(\widehat A\) maps naturally to
  \((A/q)/(\mathfrak m(A/q))^n\).
\end{definition}

\begin{proof}
  First evaluate \(\widehat A\) modulo \(\mathfrak m^n\), then pass from
  \(A/\mathfrak m^n\) to \(A/(q+\mathfrak m^n)\), and finally use the
  finite-level quotient equivalence.
\end{proof}

\begin{lemma}[Finite components on representatives]
  \label{lem:completed_ring_to_quotient_finite_of}
  \lean{Run202608192034.TODO.completedRingToQuotientFinite_of}
  \uses{def:completed_ring_to_quotient_finite}
  The finite component sends the image of \(x\in A\) in \(\widehat A\) to the
  class of \(x\) in \((A/q)/(\mathfrak m(A/q))^n\).
\end{lemma}

\begin{proof}
  This follows from the representative formula for `AdicCompletion.evalₐ` and
  the finite-level quotient equivalence.
\end{proof}

\begin{lemma}[Finite comparison on elements from \(A\)]
  \label{lem:quotient_module_level_to_ring_level_equiv_eval_map_of}
  \lean{Run202608192034.TODO.quotientModuleLevelToRingLevelEquiv_eval_map_of}
  \uses{lem:completed_ring_to_quotient_finite_of,lem:quotient_module_level_equiv_mk}
  For an element \(x\in A\), evaluating the completed module map and then
  applying the finite comparison gives the same finite quotient class as
  `completedRingToQuotientFinite`.
\end{lemma}

\begin{proof}
  Rewrite the completion map and evaluation on elements from \(A\), then use
  the representative formulas for the finite-level comparison and finite map.
\end{proof}

\begin{lemma}[Finite comparison for all completed elements]
  \label{lem:quotient_module_level_to_ring_level_equiv_eval_map}
  \lean{Run202608192034.TODO.quotientModuleLevelToRingLevelEquiv_eval_map}
  \uses{def:completed_ring_to_quotient_finite,def:quotient_module_level_to_ring_level_equiv}
  For every element of \(\widehat A\), evaluating the completed module quotient
  map and applying the finite comparison agrees with
  `completedRingToQuotientFinite`.
\end{lemma}

\begin{proof}
  Prove the identity by induction on the Cauchy-sequence representative of an
  element of the adic completion; the finite-level statement reduces to the
  representative formula.
\end{proof}

\begin{lemma}[Finite components are surjective]
  \label{lem:completed_ring_to_quotient_finite_surjective}
  \lean{Run202608192034.TODO.completedRingToQuotientFinite_surjective}
  \uses{def:completed_ring_to_quotient_finite,def:quotient_power_quotient_equiv}
  Each finite component
  \(\widehat A\to (A/q)/(\mathfrak m(A/q))^n\) is surjective.
\end{lemma}

\begin{proof}
  Pull a target element back along the finite quotient equivalence, lift it
  through the quotient map, and then use the surjectivity of evaluation from
  the adic completion at finite level.
\end{proof}

\begin{lemma}[Finite components are compatible]
  \label{lem:completed_ring_to_quotient_finite_compatible}
  \lean{Run202608192034.TODO.completedRingToQuotientFinite_compatible}
  \uses{def:completed_ring_to_quotient_finite}
  The finite maps \(\widehat A\to (A/q)/(\mathfrak m(A/q))^n\) commute with
  the transition maps in the quotient completion inverse system.
\end{lemma}

\begin{proof}
  It suffices to check a Cauchy sequence representative of an element of
  \(\widehat A\).  The Cauchy condition says the \(m\)-th and \(n\)-th terms
  differ by an element of \(\mathfrak m^m\); mapping to \(A/q\) puts the
  difference in \((\mathfrak m(A/q))^m\).
\end{proof}

\begin{definition}[The natural map to the quotient completion]
  \label{def:completed_ring_to_quotient_completion}
  \lean{Run202608192034.TODO.completedRingToQuotientCompletion}
  \uses{lem:completed_ring_to_quotient_finite_compatible}
  The compatible finite components induce a natural ring homomorphism
  \[
    \widehat A\to \widehat{A/q}.
  \]
\end{definition}

\begin{proof}
  Apply `AdicCompletion.liftRingHom` to the compatible family of finite-level
  ring maps.
\end{proof}

\begin{lemma}[Evaluating the natural map to the quotient completion]
  \label{lem:completed_ring_to_quotient_completion_eval}
  \lean{Run202608192034.TODO.completedRingToQuotientCompletion_eval}
  \uses{def:completed_ring_to_quotient_completion}
  Evaluation of the natural map \(\widehat A\to\widehat{A/q}\) at level \(n\)
  recovers the \(n\)-th finite component.
\end{lemma}

\begin{proof}
  This is the evaluation formula for `AdicCompletion.liftRingHom`.
\end{proof}

\begin{lemma}[The natural map on elements from \(A\)]
  \label{lem:completed_ring_to_quotient_completion_of_eval}
  \lean{Run202608192034.TODO.completedRingToQuotientCompletion_of_eval}
  \uses{lem:completed_ring_to_quotient_completion_eval,lem:completed_ring_to_quotient_finite_of}
  If \(x\in A\), then the level-\(n\) evaluation of its image in
  \(\widehat{A/q}\) is the class of \(x\) modulo \(q\) and
  \((\mathfrak m(A/q))^n\).
\end{lemma}

\begin{proof}
  Combine the evaluation formula with the finite-level representative formula.
\end{proof}

\begin{lemma}[The two completion maps agree]
  \label{lem:quotient_module_completion_to_quotient_adic_completion_comp_map}
  \lean{Run202608192034.TODO.quotientModuleCompletionToQuotientAdicCompletion_comp_map}
  \uses{lem:quotient_module_level_to_adic_level_equiv_eq_factor_ring_level,lem:quotient_module_level_to_ring_level_equiv_eval_map,lem:completed_ring_to_quotient_completion_eval}
  The comparison map from module completion to intrinsic quotient completion,
  composed with the completed quotient map, is the natural map
  \(\widehat A\to\widehat{A/q}\).
\end{lemma}

\begin{proof}
  Use adic-completion extensionality.  At level \(n\), rewrite the adic-level
  comparison through the ring-level factor map, use the finite comparison
  identity, and identify the result with the evaluation of
  `completedRingToQuotientCompletion`.
\end{proof}

\begin{lemma}[The natural map to the quotient completion is surjective]
  \label{lem:completed_ring_to_quotient_completion_surjective}
  \lean{Run202608192034.TODO.completedRingToQuotientCompletion_surjective}
  \uses{lem:quotient_completion_map_surjective,lem:quotient_module_completion_to_quotient_adic_completion_surjective,lem:quotient_module_completion_to_quotient_adic_completion_comp_map}
  The natural map \(\widehat A\to\widehat{A/q}\) is surjective.
\end{lemma}

\begin{proof}
  Given a point of the intrinsic quotient completion, lift it through the
  comparison equivalence to the module completion, then lift that through the
  completed quotient map \( \widehat A\to\widehat{A/q}\) at the module level.
\end{proof}

\begin{lemma}[Exactness after completing the quotient sequence]
  \label{lem:quotient_completion_module_exact}
  \lean{Run202608192034.TODO.quotient_completion_module_exact}
  \uses{lem:quotient_mk_linear_surjective,def:quotient_linear_map}
  Completing the exact sequence \(q\to A\to A/q\) remains exact at the completed
  middle term.
\end{lemma}

\begin{proof}
  Use `LinearMap.exact_subtype_mkQ` for the exactness of
  \(q\to A\to A/q\), then apply Mathlib's
  `AdicCompletion.map_exact` over the Noetherian ring \(A\).
\end{proof}

\begin{lemma}[The kernel lies in the completed submodule image]
  \label{lem:completed_ring_to_quotient_completion_ker_le_completed_submodule_range}
  \lean{Run202608192034.TODO.completedRingToQuotientCompletion_ker_le_completedSubmoduleRange}
  \uses{lem:quotient_completion_module_exact,lem:quotient_module_completion_to_quotient_adic_completion_comp_map,lem:quotient_module_completion_to_quotient_adic_completion_left_inverse}
  The kernel of the parent map
  \(\widehat A\to\widehat{A/q}\) is contained in the range of the completed
  inclusion \(\widehat q\to\widehat A\).
\end{lemma}

\begin{proof}
  If an element maps to zero in the intrinsic quotient completion, use the
  comparison equivalence to see that it maps to zero under the completed
  \(A\)-module quotient map.  Exactness of the completed sequence then puts it
  in the range of the completed inclusion of \(q\).
\end{proof}

\begin{lemma}[\(q\widehat A\) is in the kernel]
  \label{lem:completed_ring_to_quotient_completion_kernel_contains_completed_q}
  \lean{Run202608192034.TODO.completedRingToQuotientCompletion_ker_contains_completed_q}
  \uses{lem:completed_ring_to_quotient_completion_of_eval}
  The extended ideal \(q\widehat A\) is contained in the kernel of
  \(\widehat A\to\widehat{A/q}\).
\end{lemma}

\begin{proof}
  Check the claim on generators coming from \(q\).  At every finite level their
  image in \(A/q\) is zero, hence the element is zero in the quotient
  completion by extensionality of adic completions.
\end{proof}

\begin{definition}[The descended map from the completed quotient]
  \label{def:completed_quotient_to_quotient_completion}
  \lean{Run202608192034.TODO.completedQuotientToQuotientCompletion}
  \uses{lem:completed_ring_to_quotient_completion_kernel_contains_completed_q}
  The containment \(q\widehat A\subseteq\ker(\widehat A\to\widehat{A/q})\)
  gives a descended map
  \[
    \widehat A/q\widehat A\to\widehat{A/q}.
  \]
\end{definition}

\begin{proof}
  Apply the universal property of quotient rings.
\end{proof}

\begin{lemma}[The descended map on representatives]
  \label{lem:completed_quotient_to_quotient_completion_mk}
  \lean{Run202608192034.TODO.completedQuotientToQuotientCompletion_mk}
  \uses{def:completed_quotient_to_quotient_completion}
  On a representative \(x\in\widehat A\), the descended map is the natural map
  \(\widehat A\to\widehat{A/q}\).
\end{lemma}

\begin{proof}
  This follows by unfolding the quotient lift.
\end{proof}

\begin{definition}[Source facts for the completion-quotient map]
  \label{def:completed_quotient_map_source_facts}
  \lean{Run202608192034.TODO.CompletedQuotientMapSourceFacts}
  \uses{def:completed_ring_to_quotient_completion,lem:completed_ring_to_quotient_completion_kernel_contains_completed_q}
  Record the two exactness facts needed for completion commuting with quotient:
  the natural map \(\widehat A\to\widehat{A/q}\) is surjective, and its kernel
  is contained in \(q\widehat A\).
\end{definition}

\begin{proof}
  This structure isolates the two nontrivial halves of the exactness theorem
  before they are used to prove bijectivity of the descended map.
\end{proof}

\begin{lemma}[Completion-quotient source facts]
  \label{lem:completed_quotient_map_source_facts_source}
  \lean{Run202608192034.TODO.completedQuotientMapSourceFacts_source}
  \uses{def:completed_quotient_map_source_facts,lem:completed_ring_to_quotient_completion_surjective,lem:completed_ring_to_quotient_completion_ker_le_completed_q_source}
  Package the now-checked surjectivity of
  \(\widehat A\to\widehat{A/q}\) with the remaining source kernel
  containment.
\end{lemma}

\begin{proof}
  Fill the surjectivity field using
  \cref{lem:completed_ring_to_quotient_completion_surjective}; the only
  remaining source input is the reverse kernel containment.
\end{proof}

\begin{lemma}[The natural completion-to-quotient-completion map is surjective]
  \label{lem:completed_ring_to_quotient_completion_surjective_source}
  \lean{Run202608192034.TODO.completedRingToQuotientCompletion_surjective_source}
  \uses{lem:completed_ring_to_quotient_completion_surjective}
  The natural map \(\widehat A\to\widehat{A/q}\) is surjective.
\end{lemma}

\begin{proof}
  This is the checked surjectivity theorem, kept under the earlier source-name
  interface for downstream compatibility.
\end{proof}

\begin{lemma}[The descended map is surjective]
  \label{lem:completed_quotient_to_quotient_completion_surjective}
  \lean{Run202608192034.TODO.completedQuotientToQuotientCompletion_surjective}
  \uses{def:completed_quotient_to_quotient_completion,lem:completed_ring_to_quotient_completion_surjective_source}
  The descended map
  \[
    \widehat A/q\widehat A\to\widehat{A/q}
  \]
  is surjective.
\end{lemma}

\begin{proof}
  Lift a preimage along \(\widehat A\to\widehat{A/q}\), then pass it to the
  quotient \(\widehat A/q\widehat A\).
\end{proof}

\begin{lemma}[The tensor image lies in the extended ideal]
  \label{lem:completed_ideal_tensor_image_le_completed_q}
  \lean{Run202608192034.TODO.completedIdealTensorImage_le_completed_q}
  \uses{def:quotient_linear_map}
  The image of the tensor presentation
  \(\widehat A\otimes_A q\to \widehat A\) lands in the extended ideal
  \(q\widehat A\).
\end{lemma}

\begin{proof}
  Use tensor induction.  The pure tensor case is
  \(r\otimes x\mapsto r\cdot x\), which lies in the ideal generated by the
  image of \(q\); zero and addition follow from the ideal laws and linearity.
\end{proof}

\begin{lemma}[The completed submodule image lies in the extended ideal]
  \label{lem:completed_submodule_range_le_completed_q}
  \lean{Run202608192034.TODO.completedSubmoduleRange_le_completed_q}
  \uses{lem:completed_ideal_tensor_image_le_completed_q}
  The range of the completed inclusion
  \(\widehat q\to\widehat A\) is contained in \(q\widehat A\).
\end{lemma}

\begin{proof}
  Since \(q\) is a finitely generated \(A\)-module over the Noetherian ring
  \(A\), Mathlib's tensor presentation of adic completion is surjective on
  \(\widehat q\).  Naturality of that tensor presentation reduces the claim to
  \cref{lem:completed_ideal_tensor_image_le_completed_q}.
\end{proof}

\begin{lemma}[The reverse kernel containment]
  \label{lem:completed_ring_to_quotient_completion_ker_le_completed_q_source}
  \lean{Run202608192034.TODO.completedRingToQuotientCompletion_ker_le_completed_q_source}
  \uses{lem:completed_ring_to_quotient_completion_ker_le_completed_submodule_range,lem:completed_submodule_range_le_completed_q}
  The kernel of \(\widehat A\to\widehat{A/q}\) is contained in
  \(q\widehat A\).
\end{lemma}

\begin{proof}
  Exactness gives containment in the range of the completed inclusion
  \(\widehat q\to\widehat A\), and
  \cref{lem:completed_submodule_range_le_completed_q} identifies that range
  with a subideal of the algebraic extended ideal \(q\widehat A\).
\end{proof}

\begin{lemma}[The natural map has kernel exactly \(q\widehat A\)]
  \label{lem:completed_ring_to_quotient_completion_ker_eq_completed_q_source}
  \lean{Run202608192034.TODO.completedRingToQuotientCompletion_ker_eq_completed_q_source}
  \uses{lem:completed_ring_to_quotient_completion_kernel_contains_completed_q,lem:completed_ring_to_quotient_completion_ker_le_completed_q_source}
  The kernel of \(\widehat A\to\widehat{A/q}\) is exactly
  \(q\widehat A\).
\end{lemma}

\begin{proof}
  Combine the previously checked containment \(q\widehat A\subseteq\ker\)
  with the reverse source containment.
\end{proof}

\begin{lemma}[The descended map is injective]
  \label{lem:completed_quotient_to_quotient_completion_injective}
  \lean{Run202608192034.TODO.completedQuotientToQuotientCompletion_injective}
  \uses{def:completed_quotient_to_quotient_completion,lem:completed_ring_to_quotient_completion_ker_le_completed_q_source}
  The descended map
  \[
    \widehat A/q\widehat A\to\widehat{A/q}
  \]
  is injective.
\end{lemma}

\begin{proof}
  Apply the quotient-lift injectivity criterion: a lift out of
  \(\widehat A/q\widehat A\) is injective once the kernel of the parent map is
  contained in \(q\widehat A\).
\end{proof}

\begin{lemma}[The descended quotient-completion map is bijective]
  \label{lem:completed_quotient_to_quotient_completion_bijective_source}
  \lean{Run202608192034.TODO.completedQuotientToQuotientCompletion_bijective_source}
  \uses{lem:completed_quotient_to_quotient_completion_surjective,lem:completed_quotient_to_quotient_completion_injective}
  The natural map
  \[
    \widehat A/q\widehat A\to\widehat{A/q}
  \]
  is bijective.
\end{lemma}

\begin{proof}
  Combine the separately proved surjectivity and injectivity of the descended
  map.
\end{proof}

\begin{definition}[The quotient-completion equivalence]
  \label{def:quotient_completion_equiv_source}
  \lean{Run202608192034.TODO.quotientCompletionEquiv_source}
  \uses{lem:completed_quotient_to_quotient_completion_bijective_source}
  The first field of the bridge is the ring equivalence
  \[
    \widehat{A/q}\cong \widehat A/q\widehat A
  \]
  obtained from the bijective descended map.
\end{definition}

\begin{proof}
  Apply `RingEquiv.ofBijective` to the descended map and reverse the resulting
  equivalence.
\end{proof}

\begin{definition}[Quotient completion analytic bridge]
  \label{def:quotient_completion_analytic_bridge}
  \lean{Run202608192034.TODO.QuotientCompletionAnalyticBridge}
  \uses{def:quotient_adic_completion,def:quotient_completion_equiv_source}
  The remaining quotient-completion source package has two fields: the
  ring equivalence
  \[
    \widehat{A/q}\cong \widehat A/q\widehat A
  \]
  and the one-dimensional analytic criterion for \(A/q\) stated using the
  actual quotient completion.
\end{definition}

\begin{proof}
  This is a structure definition which records the two source facts separately.
\end{proof}

\begin{lemma}[Dimension-one criterion on the actual quotient completion]
  \label{lem:quotient_dimension_one_criterion_on_quotient_completion_source}
  \lean{Run202608192034.TODO.quotient_dimensionOneCriterionOnQuotientCompletion_source}
  \uses{lem:quotient_dimension_one_criterion_on_quotient_completion_of_source,lem:dimension_one_weak_criterion_adic_completion_source,lem:bad_quotient_source_data_quotient_adic_ideal_eq_maximal_ideal_source}
  The one-dimensional criterion applies directly to \(A/q\) with its actual
  adic completion \(\widehat{A/q}\).
\end{lemma}

\begin{proof}
  The dimension-one calculation is supplied by
  \cref{lem:bad_quotient_source_data_quotient_ring_krull_dim_eq_one_source}.
  The quotient Noetherian, local, and domain instances were already checked.
  The quotient ideal is identified with the quotient maximal ideal before
  applying the Basic-level criterion.
  Apply the general Anderson--Farley criterion for one-dimensional adic
  completions to \(A/q\).
\end{proof}

\begin{lemma}[Quotient completion analytic bridge source]
  \label{lem:quotient_completion_analytic_bridge_source}
  \lean{Run202608192034.TODO.quotientCompletionAnalyticBridge_source}
  \uses{def:bad_quotient_source_data,def:quotient_completion_analytic_bridge,def:quotient_completion_equiv_source,lem:quotient_dimension_one_criterion_on_quotient_completion_source}
  The paper's completion-quotient step and one-dimensional quotient criterion
  provide the bridge package for \(A/q\).
\end{lemma}

\begin{proof}
  Combine the separately named quotient-completion equivalence and the
  one-dimensional criterion on the actual quotient completion.
\end{proof}

\begin{lemma}[Analytic criterion on the completed quotient target]
  \label{lem:quotient_analytic_criterion_on_completed_quotient_source}
  \lean{Run202608192034.TODO.quotient_analyticCriterionOnCompletedQuotient_source}
  \uses{def:bad_quotient_source_data,lem:quotient_completion_analytic_bridge_source,def:completed_quotient_target_equiv_node}
  For the one-dimensional quotient \(A/q\), weak quasi-completeness is
  equivalent to analytic irreducibility of the completion target
  \(\widehat A/q\widehat A\).
\end{lemma}

\begin{proof}
  Apply the quotient completion analytic bridge, then transport the domain
  condition across the stored ring equivalence.
\end{proof}

\begin{lemma}[All quotient rings weak criterion source]
  \label{lem:all_quotients_weak_criterion_source}
  \lean{Run202608192034.TODO.all_quotients_weak_criterion_source}
  \uses{def:bad_quotient_source_data,lem:quotient_weak_characterization}
  The selected ring \(A\) satisfies the cited criterion that quasi-completeness
  is equivalent to weak quasi-completeness of all quotient rings.
\end{lemma}

\begin{proof}
  This is now a direct specialization of the general quotient characterization
  proved from the definitions in
  \cref{lem:quotient_weak_characterization}.
\end{proof}

\begin{lemma}[Transporting the dimension-one criterion]
  \label{lem:quotient_completion_dimension_criterion}
  \lean{Run202608192034.TODO.QuotientCompletionWitness.dimensionCriterion}
  \uses{def:quotient_completion_witness,def:bad_quotient_dimension_criterion_field}
  If the quotient completion target is equivalent to \(T/aT\), then the
  analytic irreducibility criterion may be stated directly with \(T/aT\) as
  the target.
\end{lemma}

\begin{proof}
  In this formalization analytic irreducibility is represented by the target
  being a domain.  Domainness is transported across the recorded ring
  equivalence.
\end{proof}

\begin{definition}[Structured bad-quotient source]
  \label{def:bad_quotient_structured_source}
  \lean{Run202608192034.TODO.BadQuotientStructuredSource}
  \uses{def:bad_quotient_source_data,def:jensen_completion_witness,def:quotient_completion_witness,def:bad_quotient_quasi_criterion_field,def:adic_completion_prime_contraction_condition}
  Package the bad quotient data together with Jensen's completion witness, the
  quotient-completion bridge, the quotient criterion for quasi-completeness,
  and the selected source ring's adic-completion prime-contraction condition.
\end{definition}

\begin{proof}
  This structure separates the paper's order into three layers: construct the
  ring \(A\), identify its completion with \(T\), and then identify the
  completion of \(A/aA\) with \(T/aT\).  It now also records the
  project-specific Farley first-direction output for the selected source ring.
\end{proof}

\begin{lemma}[Source data recovers the contracted prime]
  \label{lem:source_data_to_contracted_prime}
  \lean{Run202608192034.TODO.BadQuotientSourceData.to_contractedPrime}
  \uses{def:bad_quotient_source_data,def:contracted_prime}
  The structured source data remembers the contracted prime
  \(q=Q\cap A\), so it recovers the contracted-prime node.
\end{lemma}

\begin{proof}
  Return the stored counterexample-ring witness, the equality
  \(q=Q\cap A\), primality of \(q\), and nonzeroness of \(q\).
\end{proof}

\begin{lemma}[Source data recovers the prime generator]
  \label{lem:source_data_to_prime_generator}
  \lean{Run202608192034.TODO.BadQuotientSourceData.to_primeGenerator}
  \uses{def:bad_quotient_source_data,def:prime_generator}
  The structured source data remembers the generator \(q=aA\), so it recovers
  the prime-generator node used in the proof that \(aT\) is not prime.
\end{lemma}

\begin{proof}
  Return the stored Jensen data, weak quasi-completeness, nonzero-contraction
  condition, contracted-prime identity, and principal generator identity.
\end{proof}

\begin{lemma}[Bad-quotient source data from Jensen]
  \label{lem:bad_quotient_source_data_from_jensen}
  \lean{Run202608192034.TODO.badQuotient_sourceData_from_jensen}
  \uses{lem:prime_generator_source,def:bad_quotient_source_data}
  The prime-generator source theorem supplies the source data
  \(A,\mathfrak m_A,q,a\) used for the bad quotient.
\end{lemma}

\begin{proof}
  Destructure the prime-generator package and return the same data as a
  `BadQuotientSourceData` object.
\end{proof}

\begin{lemma}[Jensen completion equivalence source]
  \label{lem:jensen_completion_witness_source}
  \lean{Run202608192034.TODO.jensenCompletionWitness_source}
  \uses{def:jensen_completion_witness}
  The Jensen construction identifies the completion of the selected ring
  \(A\) with the prescribed complete local ring \(T\).
\end{lemma}

\begin{proof}
  The witness is now supplied explicitly by the selected source package or by
  downstream data structures; this lemma simply returns that witness.
\end{proof}

\begin{lemma}[Quotient completion source]
  \label{lem:quotient_completion_witness_source}
  \lean{Run202608192034.TODO.quotientCompletionWitness_source}
  \uses{lem:jensen_completion_witness_source,def:quotient_completion_witness,def:completed_quotient_target_equiv_node,lem:quotient_analytic_criterion_on_completed_quotient_source}
  Completion commutes with quotient by \(q=aA\), giving a completion target for
  \(A/q\) equivalent to \(T/aT\).
\end{lemma}

\begin{proof}
  The formal proof now takes the target to be
  \(\widehat A/q\widehat A\), uses the checked quotient equivalence to
  \(T/aT\), and stores the remaining one-dimensional analytic criterion as a
  separate source lemma.
\end{proof}

\begin{lemma}[Bad-quotient quasi-completeness criterion source]
  \label{lem:bad_quotient_quasi_criterion_source}
  \lean{Run202608192034.TODO.badQuotient_quasiCriterion_source}
  \uses{lem:bad_quotient_source_data_from_jensen,lem:all_quotients_weak_criterion_source}
  The selected ring \(A\) satisfies the Anderson/Farley criterion:
  \(A\) is quasi-complete exactly when every quotient \(A/J\) is weakly
  quasi-complete.
\end{lemma}

\begin{proof}
  Apply the project-level wrapper for the cited quotient characterization to
  the specialized criterion.
\end{proof}

\begin{lemma}[Source theorem for the structured bad quotient]
  \label{lem:bad_quotient_structured_source}
  \lean{Run202608192034.TODO.badQuotient_structured_source}
  \uses{def:bad_quotient_structured_source,lem:bad_quotient_source_data_from_jensen,lem:bad_quotient_source_data_jensen_completion_witness,lem:quotient_completion_witness_source,lem:bad_quotient_quasi_criterion_source,lem:bad_quotient_source_data_weakly_quasi_complete_to_adic_completion_prime_contraction_source,lem:source_data_to_contracted_prime,lem:source_data_to_prime_generator}
  The source construction supplies the complete structured package for the bad
  quotient.
\end{lemma}

\begin{proof}
  Choose the source data supplied by Jensen and the UFD generator step.  Attach
  the Jensen completion witness, the quotient-completion witness, and the
  quotient criterion.  The selected source prime-contraction condition is
  supplied by the project-specific Farley first-direction theorem.  The
  resulting tuple is the structured source package.
\end{proof}

\begin{lemma}[Source data gives a bad quotient]
  \label{lem:source_data_to_bad_quotient}
  \lean{Run202608192034.TODO.BadQuotientSourceData.to_badQuotient}
  \uses{def:bad_quotient_source_data,def:bad_quotient}
  Forgetting the two final criterion fields, the structured source data is
  exactly the bad quotient data.
\end{lemma}

\begin{proof}
  Return the fields of the source package in the order required by
  \cref{def:bad_quotient}.
\end{proof}

\begin{lemma}[Dimension and domain properties of the bad quotient]
  \label{lem:bad_quotient_dimension_domain}
  \lean{Run202608192034.TODO.badQuotient_dimension_domain}
  \uses{def:bad_quotient,def:prime_generator,lem:contracted_prime_height,lem:counterexample_ring_properties}
  The ring \(B\) is a one-dimensional Noetherian local domain.
\end{lemma}

\begin{proof}
  The ideal \(aA=q\) is prime, so the quotient is a domain; quotients of
  Noetherian local rings are Noetherian and local.  The chain
  \(q\subsetneq\mathfrak m_A\) gives dimension at least one.  A longer prime
  chain above \(q\), preceded by \((0)\subsetneq q\), would have length above
  two in \(A\), contrary to \(\dim A=2\).  Hence \(\dim B=1\).
\end{proof}

\begin{lemma}[A quotient by a nonprime ideal is not a domain]
  \label{lem:quotient_not_domain_of_not_prime}
  \lean{Run202608192034.TODO.quotient_not_domain_of_not_prime}
  If an ideal \(I\) is not prime, then the quotient \(R/I\) is not a domain.
\end{lemma}

\begin{proof}
  If \(R/I\) were a domain, then the kernel of the quotient map
  \(R\to R/I\) would be prime.  This kernel is exactly \(I\), contradiction.
\end{proof}

\begin{lemma}[Structured source criteria for the bad quotient]
  \label{lem:bad_quotient_structured_criteria_source}
  \lean{Run202608192034.TODO.badQuotient_structured_criteria_source}
  \uses{lem:bad_quotient_structured_source,lem:quotient_completion_dimension_criterion}
  The constructed data can be packaged as a source-ordered tuple containing
  \(A\), \(\mathfrak m_A\), the completion embedding \(A\to T\), the contracted
  prime \(q=Q\cap A\), a generator \(q=aA\), and the two criteria needed for
  the final bad quotient.
\end{lemma}

\begin{proof}
  Destructure the structured source theorem.  The quasi-completeness criterion
  is stored directly, and the dimension-one criterion is transported from the
  quotient completion target across its equivalence with \(T/aT\).
\end{proof}

\begin{lemma}[Expanded source criteria for the bad quotient]
  \label{lem:bad_quotient_criteria_source}
  \lean{Run202608192034.TODO.badQuotient_criteria_source}
  \uses{lem:bad_quotient_structured_criteria_source}
  Expanding the structured source package gives the explicit existential
  statement used by the downstream Lean proof.
\end{lemma}

\begin{proof}
  Destructure the source package and return its fields in the order expected by
  the public bad-quotient interface.
\end{proof}

\begin{lemma}[Completion source data for the bad quotient]
  \label{lem:bad_quotient_completion_source}
  \lean{Run202608192034.TODO.badQuotient_completion_source}
  \uses{lem:bad_quotient_criteria_source}
  The bad quotient data can be chosen with an isomorphism
  \(T/aT\cong\widehat{A/aA}\), together with the source criteria needed for
  the final quasi-completeness argument.
\end{lemma}

\begin{proof}
  Take the completion target in
  \cref{lem:bad_quotient_criteria_source} to be \(T/aT\), with the identity
  isomorphism.
\end{proof}

\begin{lemma}[The completed bad quotient is not a domain]
  \label{lem:bad_quotient_completion}
  \lean{Run202608192034.TODO.badQuotient_completion_not_domain}
  \uses{lem:bad_quotient_completion_source,lem:extended_principal_specific_not_prime,lem:quotient_not_domain_of_not_prime}
  The completion of \(B\) is isomorphic to \(T/aT\), and is not a domain.
\end{lemma}

\begin{proof}
  Completion of a Noetherian local ring commutes with quotient by a finitely
  generated ideal.  Thus
  \[
    \widehat B\cong\widehat A/a\widehat A\cong T/aT.
  \]
  By \cref{lem:extended_principal_not_prime}, \(aT\) is not prime, so this
  quotient is not a domain.
\end{proof}

\begin{lemma}[The bad quotient is not weakly quasi-complete]
  \label{lem:bad_quotient_not_weak}
  \lean{Run202608192034.TODO.badQuotient_not_weaklyQuasiComplete}
  \uses{lem:bad_quotient_dimension_domain,lem:bad_quotient_completion,lem:dimension_one_weak_criterion,def:analytically_irreducible}
  The quotient \(B=A/aA\) is not weakly quasi-complete.
\end{lemma}

\begin{proof}
  By \cref{lem:bad_quotient_dimension_domain}, \(B\) is a one-dimensional
  Noetherian local domain.  By \cref{lem:bad_quotient_completion}, its
  completion is not a domain, so \(B\) is not analytically irreducible.
  Apply \cref{lem:dimension_one_weak_criterion}.
\end{proof}

\begin{lemma}[Weak candidate with a bad quotient]
  \label{lem:a_is_weak_and_has_bad_quotient}
  \lean{Run202608192034.TODO.counterexampleRing_weak_and_bad_quotient}
  \uses{lem:counterexample_ring_weak,def:prime_generator,lem:bad_quotient_dimension_domain,lem:bad_quotient_not_weak}
  The Noetherian local domain \(A\) is weakly quasi-complete, and it has a
  prime element \(a\) such that \(A/aA\) is a one-dimensional Noetherian
  local domain which is not weakly quasi-complete.
\end{lemma}

\begin{proof}
  Weak quasi-completeness is \cref{lem:counterexample_ring_weak}; the chosen
  prime element and the asserted quotient properties are
  \cref{def:prime_generator,lem:bad_quotient_dimension_domain,
  lem:bad_quotient_not_weak}.
\end{proof}

\begin{theorem}[Anderson Problem 8a]
  \label{thm:main}
  \lean{Run202608192034.TODO.andersonProblem8a}
  \uses{lem:a_is_weak_and_has_bad_quotient,lem:quotient_weak_characterization,def:quasi_complete,def:weakly_quasi_complete}
  % SOURCE: references/anderson_2014_fulltext.txt, lines 307--339; references/farley_2016_fulltext.txt, lines 61--63.
  \textit{Source: D.~D. Anderson, Corollary~2 and the quotient criterion
  around Theorem~4, references/anderson\_2014\_fulltext.txt, lines 307--339;
  J.~D. Farley, Proposition~1,
  references/farley\_2016\_fulltext.txt, lines 61--63.}
  There exists a Noetherian local ring which is weakly quasi-complete but not
  quasi-complete.
\end{theorem}

\begin{proof}
  Take the ring \(A\) from
  \cref{lem:a_is_weak_and_has_bad_quotient}.  It is weakly quasi-complete, and
  it has a quotient \(A/aA\) which is not weakly quasi-complete.  By the
  quotient characterization \cref{lem:quotient_weak_characterization}, \(A\)
  cannot be quasi-complete.  Thus \(A\) has the required two properties.
\end{proof}
